\appendixsection{A CONNECTEDNESS PROPERTY}

Lemma 1. Let X be a connected topological space and \(\Omega\) a dense open subset of X. If, for any \(x \in X\) , there exists a neighbourhood V of x such that \(V \cap \Omega\) is connected, then \(\Omega\) is connected.

Indeed, let \(\Omega_0\) be a non-empty open and closed subset of \(\Omega\) . Let \(x \in X\) and let \(V\) be a neighbourhood of \(x\) such that \(V \cap \Omega\) is connected. If \(x \in \overline{\Omega}_0\) ,

\[
(\mathrm{V} \cap \Omega) \cap \Omega_ {0} = \mathrm{V} \cap \Omega_ {0} \neq \varnothing ,
\]

so \(V \cap \Omega \subset \Omega_{0}\) . Thus, since \(\Omega\) is dense in X, \(\overline{\Omega}_{0}\) is a neighbourhood of x. Consequently, \(\overline{\Omega}_{0}\) is non-empty, open and closed, and since X is connected, \(\overline{\Omega}_{0} = X\) . Since \(\Omega_{0}\) is closed in \(\Omega\) , this implies that \(\Omega_{0} = \Omega \cap \overline{\Omega}_{0} = \Omega\) , which proves that \(\Omega\) is connected.

Lemma 2. Let U be an open ball in \(\mathbf{C}^n\) and \(f: \mathrm{U} \to \mathbf{C}\) a holomorphic function, not identically zero. Let A be a subset of U such that \(f = 0\) on A. Then U-A is dense in U and connected.

The density of U-A follows from Differentiable and Analytic Manifolds, Results, 3.2.5. Assume first that \(n = 1\) . If \(a \in \mathrm{A}\) , the power series expansion of \(f\) about \(a\) (Differentiable and Analytic Manifolds, Results, 3.2.1) is not reduced to 0, and it follows that there exists a neighbourhood \(\mathrm{V}_a\) of \(a\) in U such that \(f\) does not vanish on \(\mathrm{V}_a - \{a\}\) . Thus, \(a\) is isolated in A, which proves that A is a discrete subset of U, hence countable since U is countable at infinity. Let \(x, y \in \mathrm{U} - \mathrm{A}\) . The union of the real affine lines joining \(x\) (resp. \(y\) ) to a point of A is meagre (General Topology, Chap. IX, §5, no. 2, Def. 2). Hence, there exists \(z \in \mathrm{U} - \mathrm{A}\) such that neither of the segments \([x, z]\) and \([y, z]\) meets A. The points \(x, y, z\) thus belong to the same connected component of U-A, which proves the lemma in the case \(n = 1\) . We turn to the general case. We can assume that A is the set of zeros of \(f\) (General Topology, Chap. I, §11, no. 1, Prop. 1). Let \(x, y \in \mathrm{U} - \mathrm{A}\) and let L be an affine line containing \(x\) and \(y\) . The restriction of \(f\) to \(\mathrm{L} \cap \mathrm{U}\) is not identically zero since \(x \in \mathrm{L} \cap \mathrm{U}\) . By what has already been proved, \(x\) and \(y\) belong to the same connected component of \((\mathrm{L} \cap \mathrm{U}) - (\mathrm{L} \cap \mathrm{A})\) and hence to the same connected component of U-A.

Lemma 3. Let X be a finite dimensional connected complex-analytic manifold and let A be a subset of X satisfying the following condition:

For any \(x \in X\) , there exists an analytic function germ \(f_{x}\) , not vanishing at x, such that the germ of A at x is contained in the germ at x of the set of zeros of \(f_{x}\) .

Then X-A is dense in X and connected.

The density of X-A follows from Differentiable and Analytic Manifolds, Results, 3.2.5. We can assume that A is closed (General Topology, Chap. I, §11, no. 1, Prop. 1). For any \(x \in \mathrm{X}\) , there exists an open neighbourhood \(\mathrm{V}\) of \(x\) and an isomorphism \(c\) from \(\mathrm{V}\) to an open ball in \(\mathbf{C}^n\) such that \(c(\mathrm{A} \cap \mathrm{V})\) is contained in the set of zeros of a holomorphic function not identically zero on \(c(\mathrm{V})\) . Then, by Lemma 2, \(\mathrm{V} \cap (\mathrm{X} - \mathrm{A})\) is connected. In view of Lemma 1, this proves that \(\mathrm{X} - \mathrm{A}\) is connected.

\specialchapter{EXERCISES}

All Lie algebras and modules over them are assumed to be finite dimensional over \(k\) ; from §3 onwards, \(k\) is assumed to be of characteristic zero.

\exercisesection{\S~1}

\begin{enumerate}
\def\labelenumi{\arabic{enumi})}
\item
  Assume that \(k\) has characteristic \(p > 0\) . Let \(V\) be a vector space, \(S\) a finite set. A map \(r: S \to \operatorname{End}(V)\) satisfies condition (AC) if and only if there exists a power \(q\) of \(p\) such that \([s^q, s'^q] = 0\) for all \(s, s' \in S\) . (Use Chap. I, §1, Exerc. 19, formula (1).)
\item
  Assume that \(k\) is perfect. Let \(V\) be a finite dimensional vector space, and \(u, v \in \operatorname{End}(V)\) . Let \(u_s, u_n, v_s, v_n\) be the semi-simple and nilpotent components of \(u, v\) . The following conditions are equivalent: (i) there exists an integer \(m\) such that \((\operatorname{ad} u)^m v = 0\) ; (ii) \(u_s\) and \(v\) commute. (To prove (i) \(\Longrightarrow\) (ii), reduce to the case where \(k\) is algebraically closed and use Lemma 1 (ii).)
\item
  We make the assumptions in no. 2. Assume that \(k\) is infinite and that condition (AC) is satisfied. Let \(k'\) be a perfect extension of \(k\) . Let \(\lambda : S \to k\) be such that \(V^{\lambda}(S) \neq 0\) . Put
\end{enumerate}

\[
\mathrm{V} ^ {\prime} = \mathrm{V} \otimes_ {k} k ^ {\prime}, \quad \mathrm{S} ^ {\prime} = \mathrm{S} \otimes_ {k} k ^ {\prime}.
\]

Let \(r' : S' \to \text{End}(V')\) be the linear map obtained from \(r\) by extension of scalars. There exists a unique map \(\lambda' : S' \to k'\) such that \(V^{\lambda}(S) \otimes_k k' = V'^{\lambda'}(S')\) . (Reduce to the case where \(V = V^{\lambda}(S)\) . Let \(P\) be a polynomial function on \(S\) and \(q\) a power of the characteristic exponent of \(k\) dividing \(\dim V\) , such that \(\lambda^q = P\) . Let \(P'\) be the polynomial function on \(S'\) which extends \(P\) . For each \(s' \in S'\) , there exists a \(\lambda'(s') \in k'\) such that \(\lambda'(s')^q = P'(s')\) . Show that the characteristic polynomial of \(r'(s')\) is \((X - \lambda'(s'))^{\dim V}\) .)

\begin{enumerate}
\def\labelenumi{\arabic{enumi})}
\setcounter{enumi}{3}
\tightlist
\item
  Assume that k has characteristic zero. Let \(\mathfrak{g} = \mathfrak{sl}(3, k)\) and let a be the subalgebra of g generated by a diagonal matrix with eigenvalues 1, -1, 0. Show that a is reductive in g, that the commutant m of a in g consists of the diagonal matrices of trace zero, and that the commutant of m in g is equal to m, and hence is distinct from a (cf.~no. 5, Remark).
\end{enumerate}

¶ 5) Assume that k is infinite. Let g be a Lie algebra and V a g-module. If n is an integer \(\geq 0\) , denote by \(V_{n}\) the set of \(v \in V\) such that \(x^{n}v = 0\) for all \(x \in g\) .

\begin{enumerate}
\def\labelenumi{\alph{enumi})}
\tightlist
\item
  Show that, if \(v \in V_n\) , \(x, y \in \mathfrak{g}\) , then
\end{enumerate}

\[
\left(\sum_ {i = 1} ^ {n} x ^ {n - i} y x ^ {i - 1}\right) v = 0.
\]

(Use the fact that \((x + ty)^n v = 0\) for all \(t\in k\) .)

Replacing \(y\) by \([x, y]\) in this formula, deduce \(^3\) that \((x^n y - y x^n)v = 0\) , and hence that \(x^n y v = 0\) .

\begin{enumerate}
\def\labelenumi{\alph{enumi})}
\setcounter{enumi}{1}
\item
  Show that \(V_{n}\) is a g-submodule of V (use a)). In particular \(\mathrm{V}^{0}(\mathfrak{g})=\bigcup_{n}\mathrm{V}_{n}\) is a g-submodule of V.
\item
  Assume that \(k = \mathbf{R}\) or \(\mathbf{C}\) , and denote by \(G\) a simply-connected Lie group with Lie algebra \(\mathfrak{g}\) ; the action of \(\mathfrak{g}\) on \(V\) defines a law of operation of \(G\) on \(V\) (Chap. III, §6, no. 1). Show that an element \(v \in V\) belongs to \(V_n\) if and only if \((s - 1)^n v = 0\) for all \(s \in G\) ; in particular \(V^0(\mathfrak{g}) = V^1(G)\) .
\end{enumerate}

\begin{enumerate}
\def\labelenumi{\arabic{enumi})}
\setcounter{enumi}{5}
\tightlist
\item
  The notations are those of Exerc. 12 of Chap. I, §3. In particular, g is a Lie algebra, M a g-module, and \(\mathrm{H}^{p}(\mathfrak{g},\mathrm{M})=\mathrm{Z}^{p}(\mathfrak{g},\mathrm{M})/\mathrm{B}^{p}(\mathfrak{g},\mathrm{M})\) is the cohomology space of degree p of g with values in M.
\end{enumerate}

\begin{enumerate}
\def\labelenumi{\alph{enumi})}
\item
  Show that \(\mathrm{B}^{p}(\mathfrak{g},\mathrm{M})\) and \(\mathrm{Z}^{p}(\mathfrak{g},\mathrm{M})\) are stable under the natural representation \(\theta\) of g on the space of cochains \(\mathrm{C}^{p}(\mathfrak{g},\mathrm{M})\) . It follows that there is a representation of g on \(\mathrm{H}^{p}(\mathfrak{g},\mathrm{M})\) . Show that this representation is trivial (use the formula \(\theta = di + id\) , loc. cit.).
\item
  Let \(\overline{k}\) be an algebraic closure of k. Let \(x \in g\) and let \(x_{M}\) be the corresponding endomorphism of M. Let \(\lambda_{1}, \ldots, \lambda_{n}\) (resp. \(\mu_{1}, \ldots, \mu_{m}\) ) be the eigenvalues (in \(\overline{k}\) ) of \(ad_{g}x\) (resp. of \(x_{M}\) ), repeated according to their multiplicity. Show that the eigenvalues of the endomorphism \(\theta(x)\) of \(C^{p}(g, M)\) are the \(\mu_{j} - (\lambda_{i_{1}} + \cdots + \lambda_{i_{p}})\) , where \(1 \leq j \leq m\) and
\end{enumerate}

\[
1 \leq i _ {1} <   i _ {2} <   \dots <   i _ {p} \leq n.
\]

Deduce, using a), that \(\mathrm{H}^{p}(\mathfrak{g},\mathrm{M})=0\) if none of the \(\mu_{j}-(\lambda_{i_{1}}+\cdots+\lambda_{i_{p}})\) is zero.

\begin{enumerate}
\def\labelenumi{\alph{enumi})}
\setcounter{enumi}{2}
\tightlist
\item
  Assume that the representation \(\mathfrak{g} \to \operatorname{End}(\mathrm{M})\) is faithful, and that \(x\) satisfies the condition:
\end{enumerate}

\[
\mu_ {j _ {1}} + \dots + \mu_ {j _ {p}} \neq \mu_ {k _ {1}} + \dots + \mu_ {k _ {p + 1}}\tag{($S_{p}$)}
\]

for all \(j_{1},\ldots,j_{p},k_{1},\ldots,k_{p+1}\in[1,m]\) .

Show that we then have \(\mathrm{H}^p (\mathfrak{g},\mathrm{M}) = 0\) (remark that the eigenvalues \(\lambda_{i}\) of \(\mathrm{ad}_{\mathfrak{g}}x\) are of the form \(\mu_j - \mu_k\) , and apply \(b\) )).

\begin{enumerate}
\def\labelenumi{\arabic{enumi})}
\setcounter{enumi}{6}
\tightlist
\item
  Let \(\mathfrak{g}\) be a nilpotent Lie algebra and V a \(\mathfrak{g}\) -module such that \(V^0(\mathfrak{g}) = 0\) . Show that \(H^p(\mathfrak{g}, V) = 0\) for all \(p \geq 0\) . (Reduce to the case where \(V = V^\lambda(\mathfrak{g})\) , with \(\lambda \neq 0\) and choose an element \(x \in \mathfrak{g}\) such that \(\lambda(x) \neq 0\) . Apply Exerc. 6 b), remarking that the \(\lambda_i\) are all zero and the \(\mu_j\) are all equal to \(\lambda(x)\) .)
\end{enumerate}

Recover the Cor. of Prop. 9 (take p = 1). \(^{4}\)

¶ 8) Assume that k is of characteristic p \textgreater{} 0. Let g be a Lie algebra over k with basis \((e_{1},\ldots,e_{n})\) . Denote by U the enveloping algebra of g and C the centre of U. For \(i = 1,\ldots,n\) , choose a non-zero p-polynomial \(f_{i}\) , of degree \(d_{i}\) , such that \(f_{i}(\mathrm{ad}e_{i}) = 0\) ; then \(f_{i}(e_{i}) \in \mathrm{C}\) , cf.~Chap. I, §7, Exerc. 5. Put \(z_{i} = f_{i}(e_{i})\) .

\begin{enumerate}
\def\labelenumi{\alph{enumi})}
\item
  Show that \(z_{1},\ldots,z_{n}\) are algebraically independent. If \(A=k[z_{1},\ldots,z_{n}]\) , show that U is a free A-module with basis the monomials \(e_{1}^{\alpha_{1}}\ldots e_{n}^{\alpha_{n}}\) , where \(0\leq\alpha_{i}\leq d_{i}\) . (Use the Poincaré-Birkhoff-Witt theorem.) The rank {[}U:A{]} of U over A is equal to \(d_{1}\ldots d_{n}\) ; it is a power of p.~Deduce that C is an A-module of finite type, hence a k-algebra of finite type and of dimension n (Commutative Algebra, Chap. VIII).
\item
  Let \(\mathrm{K}\) be the field of fractions of \(\mathrm{A}\) , and let
\end{enumerate}

\[
\mathrm{U} _ {(\mathrm{K})} = \mathrm{U} \otimes_ {\mathrm{A}} \mathrm{K}, \quad \mathrm{C} _ {(\mathrm{K})} = \mathrm{C} \otimes_ {\mathrm{A}} \mathrm{K}.
\]

Then \(\mathrm{U}_{(\mathrm{K})}\supset\mathrm{C}_{(\mathrm{K})}\supset\mathrm{K}\) . Show that \(\mathrm{U}_{(\mathrm{K})}\) is a field with centre \(\mathrm{C}_{(\mathrm{K})}\) , and that this is the quotient field (both left and right) of U, cf.~Chap. I, §2, Exerc. 10. Deduce that \([\mathrm{U}_{(\mathrm{K})}:\mathrm{C}_{(\mathrm{K})}]\) is of the from \(q^{2}\) , where q is a power of p; we have \([\mathrm{C}_{(\mathrm{K})}:\mathrm{K}]=q_{\mathrm{C}}\) , where \(q_{C}\) is a power of p, and \([U:A]=q_{C}q^{2}\) .

\begin{enumerate}
\def\labelenumi{\alph{enumi})}
\setcounter{enumi}{2}
\tightlist
\item
  Let d be a non-zero element of A, and let \(\Lambda\) be a subring of \(\mathrm{U}_{(\mathrm{K})}\) such that \(U \subset \Lambda \subset d^{-1}U\) . Show that \(\Lambda = U\) . {[}If x = b/a, \(a \in A - \{0\}\) , is an element of \(\Lambda\) , show by induction on m that the relation \(b \in Ua + U_m\) implies that \(b \in Ua + U_{m-1}\) , where \(\{U_m\}\) is the canonical filtration of U. (For this, use the fact that gr U is integrally closed, and argue as in Prop. 15 of Commutative Algebra, Chap. V, §1, no. 4.) For m = 0, this gives \(b \in Ua\) , i.e.~\(x \in U\) .{]}
\end{enumerate}

Deduce that C is integrally closed.

\begin{enumerate}
\def\labelenumi{\alph{enumi})}
\setcounter{enumi}{3}
\tightlist
\item
  Assume that k is algebraically closed. Let \(\rho: \mathfrak{g} \to \mathfrak{gl}(V)\) be an irreducible linear representation of g and \(\rho_{U}\) the corresponding representation of U. The restriction of \(\rho_{U}\) to C is a homomorphism \(\gamma_{\rho}\) from C to k (identified with the homotheties of V); let \(\alpha_{\rho}\) be its restriction to A. Show that for any homomorphism \(\alpha\) (resp. \(\gamma\) ) from the k-algebra A (resp. C) to k, there exists at least one irreducible representation \(\rho\) of g such that \(\alpha_{\rho} = \alpha\) (resp. \(\gamma_{\rho} = \gamma\) ) and that there are only finitely-many such representations (up to equivalence). Show that \(\dim V \leq q\) , with the notation of b). \(^{5}\)
\end{enumerate}

¶ 9) We retain the notations of the preceding exercise, and assume further that g is nilpotent.

\begin{enumerate}
\def\labelenumi{\alph{enumi})}
\item
  Show that the basis \((e_{1},\ldots,e_{n})\) can be chosen so that, for any pair \((i,j)\) , \([e_{i},e_{j}]\) is a linear combination of the \(e_{h}\) for \(h > \sup(i,j)\) . Assume from now on that the \(e_{i}\) satisfy this condition. For \(i=1,\ldots,n\) choose a power \(q(i)\) of p such that \(\operatorname{ad}(e_{i})^{q(i)} = 0\) , and put \(z_{i} = e_{i}^{q(i)}, A = k[z_{1},\ldots,z_{n}]\) , cf.~Exerc. 8. b) Let \(\rho : g \to gl(V)\) be a linear representation of g. Assume that \(\rho(e_{i})\) is nilpotent for \(i = 1,\ldots,n\) . Show that \(\rho(x)\) is nilpotent for all \(x \in g\) . (Argue by induction on \(n = \dim g\) and reduce to the case where \(\rho\) is irreducible. Show that, in this case, \(\rho(e_{n}) = 0\) and apply the induction hypothesis.)
\item
  Let \(\rho_1: \mathfrak{g} \to \mathfrak{gl}(V_1)\) and \(\rho_2: \mathfrak{g} \to \mathfrak{gl}(V_2)\) be two linear representations of \(\mathfrak{g}\) . Assume that \(V_1\) and \(V_2\) are \(\neq 0\) , and that \(V_1 = V^{\lambda_1}(\mathfrak{g}), V_2 = V^{\lambda_2}(\mathfrak{g})\) , where \(\lambda_1\) and \(\lambda_2\) are two functions on \(\mathfrak{g}\) , cf.~no. 3. Show that, if \(\lambda_1(e_i) = \lambda_2(e_i)\) for \(i = 1, \ldots, n\) , then \(\lambda_1 = \lambda_2\) and there exists a non-zero \(\mathfrak{g}\) -homomorphism from \(V_1\) to \(V_2\) (apply \(b\) ) to the \(\mathfrak{g}\) -module \(V = \mathcal{L}(V_1, V_2)\) and use Engel's theorem to show that \(V\) contains a non-zero \(\mathfrak{g}\) -invariant element). Deduce that if, in addition, \(V_1\) and \(V_2\) are simple, they are isomorphic.
\item
  Assume that k is algebraically closed. Let R be the set of equivalence classes of irreducible representations of g. If \(\rho \in R\) , put
\end{enumerate}

\[
x _ {\rho} = (x _ {\rho} (1), \dots , x _ {\rho} (n)) \in k ^ {n},
\]

where \(x_{\rho}(i)\) is the unique eigenvalue of \(\rho(e_{i})\) . Show that \(\rho \mapsto x_{\rho}\) is a bijection from R to \(k^{n}\) . (Injectivity follows from c), and surjectivity from Exerc. 8 d). Deduce the following consequences:

\begin{enumerate}
\def\labelenumi{(\roman{enumi})}
\item
  For any maximal ideal \(\mathfrak{m}\) of A, the quotient of U/mU by its radical is a matrix algebra.
\item
  The degree of any irreducible representation of \(\mathfrak{g}\) is a power of \(p\) (this follows from (i) and the fact that \([\mathrm{U} / \mathfrak{mU}:k]\) is a power of \(p\) ).
\item
  Every homomorphism from A to k extends uniquely to a homomorphism from C to k (use the fact that C/mC is contained in the centre of U/mU, which is a local k-algebra with residue field k).
\item
  There exists an integer N ≥ 0 such that \(x^{p^{N}} \in A\) for all \(x \in C\) (this follows from (iii)). \(^{6}\)
\end{enumerate}

¶ 10) Assume that \(k\) is of characteristic \(p > 0\) . Denote by \(\mathfrak{g}\) a Lie algebra with basis \(\{e_1, e_2, e_3\}\) , with \([e_1, e_2] = e_3\) , \([e_1, e_3] = [e_2, e_3] = 0\) .

\begin{enumerate}
\def\labelenumi{\alph{enumi})}
\item
  Show that the centre of \(\mathrm{U}\mathfrak{g}\) is \(k[e_1^p,e_2^p,e_3]\) .
\item
  Assume that k is algebraically closed. Show that, for all \((\lambda_{1},\lambda_{2},\lambda_{3})\in k^{3}\) , there exists (up to equivalence) a unique irreducible representation \(\rho\) of g such that \(\lambda_{i}\) is the unique eigenvalue of \(\rho(e_i)\) ( \(i = 1,2,3\) ); the degree of \(\rho\) is \(p\) if \(\lambda_3 \neq 0\) , and is 1 if \(\lambda_3 = 0\) . (Apply Exerc. 8 and 9, or argue directly.)
\end{enumerate}

\begin{enumerate}
\def\labelenumi{\arabic{enumi})}
\setcounter{enumi}{10}
\tightlist
\item
  Let \(\mathfrak{h}\) be a nilpotent Lie algebra, \(V\) an \(\mathfrak{h}\) -module not reduced to 0, and \(\lambda\) a function on \(\mathfrak{h}\) such that \(V = V^{\lambda}(\mathfrak{h})\) . Prove the equivalence of the following properties:
\end{enumerate}

\begin{enumerate}
\def\labelenumi{(\roman{enumi})}
\item
  \(\lambda\) is a linear form on \(\mathfrak{h}\) , zero on \([\mathfrak{h},\mathfrak{h}]\) .
\item
  There exists a basis of \(\mathrm{V}\) with respect to which the endomorphisms defined by the elements of \(\mathfrak{h}\) are triangular.
\end{enumerate}

(To prove that (i) \(\Longrightarrow\) (ii), apply Engel's theorem to the \(\mathfrak{h}\) -module \(\mathcal{L}(\mathrm{W}_{\lambda},\mathrm{V})\) , where \(\mathrm{W}\) is the 1-dimensional \(\mathfrak{h}\) -module defined by \(\lambda\) .)

Properties (i) and (ii) are true if k is of characteristic 0 (Prop. 9).

\exercisesection{\S~2}

\begin{enumerate}
\def\labelenumi{\arabic{enumi})}
\item
  The diagonal matrices of trace 0 form a Cartan subalgebra of \(\mathfrak{sl}(n,k)\) , except when \(n = 2\) and \(k\) is of characteristic 2.
\item
  Let \(e\) be the element \(\begin{pmatrix} 0 & 1 \\ 0 & 0 \end{pmatrix}\) of \(\mathfrak{sl}(2, \mathbf{C})\) . Show that \(\mathbf{C}e\) is a maximal nilpotent Lie subalgebra of \(\mathfrak{sl}(2, \mathbf{C})\) , but not a Cartan subalgebra of \(\mathfrak{sl}(2, \mathbf{C})\) .
\item
  Assume that \(k\) is of characteristic 0. Let \(\mathfrak{g}\) be a semi-simple Lie algebra. Let \(E\) be the set of commutative subalgebras of \(\mathfrak{g}\) all of whose elements are semi-simple in \(\mathfrak{g}\) . Then the Cartan subalgebras of \(\mathfrak{g}\) are the maximal elements of \(E\) . (Use Th. 2 and Prop. 10.)
\end{enumerate}

In particular, the union of the Cartan subalgebras of \(\mathfrak{g}\) is equal to the set of semi-simple elements of \(\mathfrak{g}\) .

\begin{enumerate}
\def\labelenumi{\arabic{enumi})}
\setcounter{enumi}{3}
\item
  Let \(\mathfrak{g}\) be a Lie algebra with a basis \((x,y,z)\) such that \([x,y] = y\) , \([x,z] = z\) , \([y,z] = 0\) . Let \(\mathfrak{a}\) be the ideal \(ky + kz\) of \(\mathfrak{g}\) . Then \(\mathrm{rk}(\mathfrak{a}) = 2\) and \(\mathrm{rk}(\mathfrak{g}) = 1\) .
\item
  Assume that \(k\) is of characteristic 0. Let \(\mathfrak{g}\) be a Lie algebra, \(\mathfrak{r}\) its radical, \(\mathfrak{h}\) a Cartan subalgebra of \(\mathfrak{g}\) . Show that
\end{enumerate}

\[
\mathfrak {r} = [ \mathfrak {g}, \mathfrak {r} ] + (\mathfrak {h} \cap \mathfrak {r}).
\]

(Observe that the image of \(\mathfrak{h}\) in \(\mathfrak{g}/[\mathfrak{g},\mathfrak{r}]\) contains the centre \(\mathfrak{r}/[\mathfrak{g},\mathfrak{r}]\) of \(\mathfrak{g}/[\mathfrak{g},\mathfrak{r}]\) .)

\begin{enumerate}
\def\labelenumi{\arabic{enumi})}
\setcounter{enumi}{5}
\item
  Let \(\mathfrak{g}\) be a Lie algebra, \(\mathfrak{h}\) a nilpotent subalgebra of \(\mathfrak{g}\) . If \(\mathfrak{g}^0 (\mathfrak{h})\) is nilpotent, \(\mathfrak{g}^0 (\mathfrak{h})\) is a Cartan subalgebra of \(\mathfrak{g}\) .
\item
  Let \(\mathfrak{s}\) be a Lie algebra, \(\mathfrak{a}\) a Cartan subalgebra of \(\mathfrak{s}\) and V an \(\mathfrak{s}\) -module. Let \(\mathfrak{g} = \mathfrak{s} \times V\) be the semi-direct product of \(\mathfrak{s}\) by V. Show that \(\mathfrak{a} \times V^0(\mathfrak{a})\) is a Cartan subalgebra of \(\mathfrak{g}\) .
\item
  Assume that \(k\) is of characteristic \(p > 0\) . Denote by \(\mathfrak{s}\) a Lie algebra with basis \(\{x, y\}\) such that \([x, y] = y\) . Let \(V\) be a \(k\) -vector space with basis \(\{e_i\}_{i \in \mathbf{Z}/p\mathbf{Z}}\) .
\end{enumerate}

\begin{enumerate}
\def\labelenumi{\alph{enumi})}
\item
  Show that V has a unique \(\mathfrak{s}\) -module structure such that \(xe_{i} = ie_{i}\) and \(ye_{i} = e_{i + 1}\) for all \(i\) . This \(\mathfrak{s}\) -module is simple.
\item
  Let \(\mathfrak{g} = \mathfrak{s} \times \mathrm{V}\) be the semi-direct product of \(\mathfrak{s}\) by \(\mathrm{V}\) . Show that \(\mathfrak{g}\) is a solvable algebra of rank 1 whose derived algebra is not nilpotent.
\item
  An element of \(\mathfrak{g}\) is regular if and only if its projection on \(\mathfrak{s}\) is of the form \(ax + by\) , with \(ab \neq 0\) .
\item
  We have \(\mathrm{V}^{0}(x+y)=0\) and \(\mathrm{V}^{0}(x)=ke_{0}\) . Deduce (cf.~Exerc. 6) that g has Cartan subalgebras of dimension 1 (for example that generated by \(x+y\) ) and Cartan subalgebras of dimension 2 (for example that generated by x and \(e_{0}\) ).
\end{enumerate}

\begin{enumerate}
\def\labelenumi{\arabic{enumi})}
\setcounter{enumi}{8}
\item
  Let g be a Lie algebra with a basis \((x,y)\) such that \([x,y]=y\) . Let k=ky, and \(\varphi:g\to g/k\) the canonical morphism. The element 0 of g/k is regular in g/k but is not the image under \(\varphi\) of a regular element of g.
\item
  Assume that \(k\) is infinite. Let \(\mathfrak{g}\) be a Lie algebra. Prove the equivalence of the following properties:
\end{enumerate}

\begin{enumerate}
\def\labelenumi{(\roman{enumi})}
\item
  \(\operatorname{rk}(\mathfrak{g}) = \dim (\mathfrak{g})\) .
\item
  \(\mathfrak{g}\) is nilpotent.
\item
  g has only finitely-many Cartan subalgebras of dimension rk(g).
\item
  g has only one Cartan subalgebra.
\end{enumerate}

\begin{enumerate}
\def\labelenumi{\arabic{enumi})}
\setcounter{enumi}{10}
\tightlist
\item
  Let \(\mathfrak{h}\) be a commutative Lie algebra \(\neq 0\) , \(P\) a finite subset of \(\mathfrak{h}^*\) containing 0. Show that there exists a Lie algebra \(\mathfrak{g}\) containing \(\mathfrak{h}\) as a Cartan subalgebra, and such that the set of weights of \(\mathfrak{h}\) in \(\mathfrak{g}\) is \(P\) . (Construct \(\mathfrak{g}\) as the semi-direct product of \(\mathfrak{h}\) by the \(\mathfrak{h}\) -module \(V\) which is the direct sum of the 1-dimensional modules corresponding to the elements of \(P - \{0\}\) , cf.~Exerc. 7.)
\end{enumerate}

An element \(x\) of \(\mathfrak{h}\) is such that \(\mathfrak{h} = \mathfrak{g}^0(x)\) if and only if \(x\) is not orthogonal to any element of \(\mathrm{P} - \{0\}\) .

\begin{enumerate}
\def\labelenumi{\arabic{enumi})}
\setcounter{enumi}{11}
\tightlist
\item
  Assume that \(k\) is finite. Construct an example of a Lie algebra \(\mathfrak{g}\) having a Cartan subalgebra \(\mathfrak{h}\) in which there exists no element \(x\) such that \(\mathfrak{h} = \mathfrak{g}^0(x)\) . (Use the preceding exercise, and take \(\mathrm{P} = \mathfrak{h}^*\) .)
\end{enumerate}

¶ 13) Assume that k is finite. Denote by \(k'\) an infinite extension of k. Let g be a Lie algebra over k. The rank of g, denoted by \(\operatorname{rk}(\mathfrak{g})\) , is the rank of the \(k'\) -Lie algebra \(g' = g \otimes_{k} k'\) ; an element of g is said to be regular if it is regular in \(g'\) ; these definitions do not depend on the choice of \(k'\) . Show that, if

\[
\operatorname{Card} (k) \geq \dim \mathfrak {g} - \operatorname{rk} (\mathfrak {g}),
\]

\(\mathfrak{g}\) contains a regular element (hence also a Cartan subalgebra).

(Use the following result: if \(a\) is a non-zero homogeneous element of \(k[\mathrm{X}_1, \ldots, \mathrm{X}_n]\) , and if \(\operatorname{Card}(k) \geq \deg(a)\) , there exists \(x \in k^n\) such that \(a(x) \neq 0\) .)

\begin{enumerate}
\def\labelenumi{\arabic{enumi})}
\setcounter{enumi}{13}
\tightlist
\item
  Assume that \(k\) is of characteristic zero. Let V be a finite dimensional \(k\) -vector space, \(\mathfrak{g}\) a Lie subalgebra of \(\mathfrak{gl}(V)\) , \(\mathfrak{h}\) a Cartan subalgebra of \(\mathfrak{g}\) and \(\mathfrak{n}_{\mathrm{V}}\) the largest ideal of nilpotency of the \(\mathfrak{g}\) -module V (Chap. I, §4, no. 3, Def. 2).
\end{enumerate}

Show that an element of h is nilpotent if and only if it belongs to \(n_{V}\) . (Reduce to the case where the g-module V is semi-simple, and use Cor. 3 of Th. 2.)

¶ 15) Assume that \(k\) is infinite. Let \(\mathfrak{g}\) be a Lie algebra, \(\mathcal{C}_{\infty}\mathfrak{g}\) the union of the ascending central series of \(\mathfrak{g}\) , and \(x\) an element of \(\mathfrak{g}\) . Prove the equivalence of the following properties:

\begin{enumerate}
\def\labelenumi{(\roman{enumi})}
\item
  \(x\) belongs to every Cartan subalgebra of \(\mathfrak{g}\) .
\item
  \(x\in \mathfrak{g}^0 (y)\) for all \(y\in \mathfrak{g}\) (i.e.~\(x\in \mathfrak{g}^0 (\mathfrak{g})\) ).
\item
  \(x \in C_{\infty}g.\)
\end{enumerate}

(The implications (iii) \(\Longrightarrow\) (ii) \(\Longrightarrow\) (i) are immediate. To prove that (i) \(\Longrightarrow\) (ii), remark that (i) is equivalent to saying that \(x\in \mathfrak{g}^0 (y)\) for every regular element \(y\) in \(\mathfrak{g}\) , and use the fact that the regular elements are dense in \(\mathfrak{g}\) in the Zariski topology. To prove that (ii) \(\Longrightarrow\) (iii), observe that \(\mathfrak{n} = \mathfrak{g}^0 (\mathfrak{g})\) is stable under \(\mathfrak{g}\) (§1, Exerc. 5) and apply Engel's theorem to the \(\mathfrak{g}\) -module \(\mathfrak{n}\) ; deduce that \(\mathfrak{n}\) is contained in \(\mathcal{L}_{\infty}\mathfrak{g}\) .)

¶ 16) Let g be a solvable complex Lie algebra, h a Cartan subalgebra of g, \(\mathfrak{g} = \bigoplus \mathfrak{g}^{\lambda}(\mathfrak{h})\) the corresponding decomposition of g into primary subspaces, with \(\mathfrak{g}^{0}(\mathfrak{h}) = \mathfrak{h}\) .

\begin{enumerate}
\def\labelenumi{\alph{enumi})}
\item
  Show that the restrictions to \(\mathfrak{h}\) of the linear forms called roots of \(\mathfrak{g}\) in Chap. III, §9, Exerc. 17 c) are the weights of \(\mathfrak{h}\) in \(\mathfrak{g}\) , i.e.~the \(\lambda\) such that \(\mathfrak{g}^{\lambda}(\mathfrak{h}) \neq 0\) ; deduce that such a \(\lambda\) vanishes on \(\mathfrak{h} \cap \mathcal{D}\mathfrak{g}\) .
\item
  Let \((x,y)\mapsto [x,y]'\) be the alternating bilinear map from \(\mathfrak{g}\times \mathfrak{g}\) to \(\mathfrak{g}\) with the following properties:
\end{enumerate}

\begin{enumerate}
\def\labelenumi{(\roman{enumi})}
\item
  If \(x \in \mathfrak{g}^{\lambda}(\mathfrak{h}), y \in \mathfrak{g}^{\mu}(\mathfrak{h})\) , with \(\lambda \neq 0, \mu \neq 0\) , then \([x, y]' = [x, y]\) ;
\item
  if \(x \in \mathfrak{g}^{0}(\mathfrak{h}), y \in \mathfrak{g}^{\mu}(\mathfrak{h})\) , then \([x, y]' = [x, y] - \mu(x)y\) .
\end{enumerate}

Show that this gives a new Lie algebra structure on \(\mathfrak{g}\) (use \(a\) ). Denote it by \(\mathfrak{g}'\) .

\begin{enumerate}
\def\labelenumi{\alph{enumi})}
\setcounter{enumi}{2}
\tightlist
\item
  Show that, if \(x \in \mathfrak{g}^{\lambda}(\mathfrak{h})\) , the map \(ad'x : y \mapsto [x, y]'\) is nilpotent. Deduce that \(g'\) is nilpotent (apply Exerc. 11 of Chap. I, §4 to the set E of \(ad'x\) , where x belongs to the union of the \(\mathfrak{g}^{\lambda}(\mathfrak{h})\) ).
\end{enumerate}

\exercisesection{\S~3}

\begin{enumerate}
\def\labelenumi{\arabic{enumi})}
\item
  Let g be a Lie algebra, \(g'\) a Cartan subalgebra of g. Then the conditions of Prop. 3 are satisfied. But an element of \(g'\) , even if it is regular in \(g'\) , is not necessarily regular in g.
\item
  Let g be a real Lie algebra of dimension n, U (resp. H) the set of regular elements (resp. of Cartan subalgebras) of g, and \(\text{Int}(\mathfrak{g})\) the group of inner automorphisms of g (Chap. III, §6, no. 2, Def. 2).
\end{enumerate}

\begin{enumerate}
\def\labelenumi{\alph{enumi})}
\item
  Show that, if \(x\) and \(y\) belong to the same connected component of \(\mathrm{U}\) , \(\mathfrak{g}^0(x)\) and \(\mathfrak{g}^0(y)\) are conjugate under \(\operatorname{Int}(\mathfrak{g})\) .
\item
  Show that the number of connected components of U is finite, and that this number is bounded by a constant \(c(n)\) depending only on n (apply Exerc. 2 of App. II).
\item
  Deduce that the number of orbits of \(\operatorname{Int}(\mathfrak{g})\) on \(\mathrm{H}\) is \(\leq c(n)\) .
\end{enumerate}

\begin{enumerate}
\def\labelenumi{\arabic{enumi})}
\setcounter{enumi}{2}
\tightlist
\item
  Let \(\mathfrak{g}\) be a real Lie algebra, \(\mathfrak{r}\) the radical of \(\mathfrak{g}\) , \(\mathfrak{h}\) and \(\mathfrak{h}'\) Cartan subalgebras of \(\mathfrak{g}\) , \(\varphi\) the canonical homomorphism from \(\mathfrak{g}\) to \(\mathfrak{g} / \mathfrak{r}\) . The following conditions are equivalent:
\end{enumerate}

\begin{enumerate}
\def\labelenumi{(\roman{enumi})}
\item
  \(\mathfrak{h}\) and \(\mathfrak{h}'\) are conjugate under \(\operatorname{Int}(\mathfrak{g})\) ;
\item
  \(\varphi(\mathfrak{h})\) and \(\varphi(\mathfrak{h}')\) are conjugate under \(\operatorname{Int}(\mathfrak{g}/\mathfrak{r})\) . (Imitate the proof of Prop. 5.)
\end{enumerate}

\begin{enumerate}
\def\labelenumi{\arabic{enumi})}
\setcounter{enumi}{3}
\item
  Let \(\mathfrak{g} = \mathfrak{sl}(2, \mathbf{R})\) , \(x = \begin{pmatrix} 1 & 0 \\ 0 & -1 \end{pmatrix}\) , \(y = \begin{pmatrix} 0 & -1 \\ 1 & 0 \end{pmatrix}\) . Show that \(\mathbf{R}x\) and \(\mathbf{R}y\) are Cartan subalgebras of \(\mathfrak{g}\) not conjugate under \(\operatorname{Aut}(\mathfrak{g})\) .
\item
  \begin{enumerate}
  \def\labelenumii{\alph{enumii})}
  \tightlist
  \item
    Show that there exists a Lie algebra \(\mathfrak{g}\) over \(k\) with basis \((x,y,z,t)\) such that
  \end{enumerate}
\end{enumerate}

\[
[ x, y ] = z, [ x, t ] = t, [ y, t ] = 0, [ \mathfrak {g}, z ] = 0.
\]

Show that g is solvable and that \(k = kx + ky + kz\) is a subalgebra of g.

\begin{enumerate}
\def\labelenumi{\alph{enumi})}
\setcounter{enumi}{1}
\item
  Show that the elementary automorphisms of g are the maps of the form \(1 + \lambda \operatorname{ad}_{\mathfrak{g}} y + \mu \operatorname{ad}_{\mathfrak{g}} t\) where \(\lambda, \mu \in k\) .
\item
  Show that \(1 + ad_{k}x\) is an elementary automorphism of k which does not extend to an elementary automorphism of g.
\end{enumerate}

\(d\) ) Let \(\mathfrak{s}\) be a semi-simple subalgebra of a Lie algebra \(\mathfrak{a}\) . Show that every elementary automorphism of \(\mathfrak{s}\) extends to an elementary automorphism of \(\mathfrak{a}\) .

\begin{enumerate}
\def\labelenumi{\arabic{enumi})}
\setcounter{enumi}{5}
\item
  Every element of a reductive Lie algebra \(\mathfrak{g}\) is contained in a commutative subalgebra of dimension \(\mathrm{rk}(\mathfrak{g})\) .
\item
  Let \(\mathfrak{g}\) be a Lie algebra and \(\mathfrak{g}'\) a subalgebra of \(\mathfrak{g}\) reductive in \(\mathfrak{g}\) . Let \(\mathfrak{a}\) be a Cartan subalgebra of \(\mathfrak{g}'\) . Show that there exists a Cartan subalgebra of \(\mathfrak{g}\) which contains \(\mathfrak{a}\) (use Prop. 10 of §2). Deduce that \(\mathrm{rk}(\mathfrak{g}') \leq \mathrm{rk}(\mathfrak{g})\) and that equality holds if and only if \(\mathfrak{g}'\) has properties (i), (ii), (iii) of Prop. 3.
\item
  Let \(\mathfrak{g}\) be a Lie algebra, \(\mathfrak{a}\) an ideal of \(\mathfrak{g}\) , \(\mathfrak{h}\) a Cartan subalgebra of \(\mathfrak{g}\) and \(\mathcal{C}_{\infty}\mathfrak{g}\) the union of the ascending central series of \(\mathfrak{g}\) . Show that \(\mathfrak{a} \subset \mathfrak{h}\) implies \(\mathfrak{a} \subset \mathcal{C}_{\infty}\mathfrak{g}\) (in other words \(\mathcal{C}_{\infty}\mathfrak{g}\) is the largest ideal of \(\mathfrak{g}\) contained in \(\mathfrak{h}\) ). (Reduce to the case where \(k\) is algebraically closed and remark that \(\mathfrak{a}\) is stable under every elementary automorphism of \(\mathfrak{g}\) ; the relation \(\mathfrak{a} \subset \mathfrak{h}\) then implies that \(\mathfrak{a}\) is contained in every Cartan subalgebra of \(\mathfrak{g}\) ; conclude by means of Exerc. 15 of §2.)
\end{enumerate}

¶ 9) Let g be a Lie algebra, h a Cartan subalgebra of g and x an element of h. Let \(g = h \oplus g^{+}\) be the Fitting decomposition (§1, no. 1) of g with respect to h.

\begin{enumerate}
\def\labelenumi{\alph{enumi})}
\item
  Let \(\mathfrak{n}\) be the largest semi-simple \(\mathfrak{h}\) -submodule contained in \(\mathfrak{g}^0 (x)\cap \mathfrak{g}^+\) . Show that \(\mathfrak{n} = 0\) if and only if \(\mathfrak{g}^0 (x) = \mathfrak{h}\) , i.e.~\(x\) is regular in \(\mathfrak{g}\) .
\item
  Show that \(\mathfrak{h} \oplus \mathfrak{n}\) is a subalgebra of \(\mathfrak{g}\) . If \(\mathfrak{h}'\) is the intersection of \(\mathfrak{h}\) with the commutant of \(\mathfrak{n}\) , show that \(\mathfrak{h}'\) is an ideal of \(\mathfrak{h} \oplus \mathfrak{n}\) which contains \(\mathcal{D}\mathfrak{h}\) and \(x\) . Conclude (Exerc. 8) that \(\mathfrak{h}' \subset \mathcal{C}_{\infty}(\mathfrak{h} \oplus \mathfrak{n})\) , and hence that \(x \in \mathcal{C}_{\infty}(\mathfrak{h} \oplus \mathfrak{n})\) and that \(x\) belongs to every Cartan subalgebra of \(\mathfrak{h} \oplus \mathfrak{n}\) .
\item
  If \(n \neq 0\) , \(h \oplus n\) is not nilpotent and has infinitely-many Cartan subalgebras (§2, Exerc. 10). Conclude that x belongs to infinitely-many Cartan subalgebras of g.
\end{enumerate}

\(d\) ) An element of \(\mathfrak{g}\) is regular if and only if it belongs to a unique Cartan subalgebra.

¶ 10) Let g be a Lie algebra, r its radical, n its largest nilpotent ideal and a one of its Levi subalgebras.

\begin{enumerate}
\def\labelenumi{\alph{enumi})}
\item
  Put \(\mathfrak{g}' = \mathfrak{n} + \mathcal{D}\mathfrak{g}\) . Show that \(\mathfrak{g}' = \mathfrak{n} \oplus \mathfrak{a}\) . (Use the fact that \([\mathfrak{g}, \mathfrak{r}]\) is contained in \(\mathfrak{n}\) .) If \(\mathfrak{g} \neq 0\) , then \(\mathfrak{g}' \neq 0\) .
\item
  Assume that k is algebraically closed. Let \((\mathrm{V}_{i})_{i\in\mathrm{I}}\) be the quotients of a Jordan-Hölder sequence of the g-module g (with the adjoint representation). If \(x\in r\) , show that \(x_{V_{i}}\) is a homothety and that \(x_{V_{i}}=0\) for all i if and only if x belongs to n.~Deduce that an element \(y\in g\) belongs to \(g'\) if and only if \(\operatorname{Tr}(y_{\mathrm{V}_{i}})=0\) for all \(i\in I\) .
\item
  Denote by N the vector subspace of g generated by the elements x such that ad x is nilpotent. Show that N is a subalgebra of g (use the fact that N is stable under \(\operatorname{Aut}_{e}(\mathfrak{g})\) ). Show, by using b), that \(N \subset g'\) .
\end{enumerate}

\(d\) ) Let \(\mathfrak{h}\) be a Cartan subalgebra of \(\mathfrak{g}\) . Assume that there exists a subset \(R\) of \(\mathfrak{h}^*\) such that

\[
\mathfrak {g} = \mathfrak {h} \oplus \bigoplus_ {\alpha \in R} \mathfrak {g} ^ {\alpha} (\mathfrak {h}),
\]

an assumption which is satisfied, in particular, if k is algebraically closed. Show that N then contains the \(\mathfrak{g}^{\alpha}(\mathfrak{h})\) , Dh and n; deduce that N contains \(g'\) , so \(N = g'\) .

\begin{enumerate}
\def\labelenumi{\alph{enumi})}
\setcounter{enumi}{4}
\tightlist
\item
  If \(k\) is algebraically closed and \(\mathfrak{g} \neq 0\) , \(\mathfrak{g}\) contains an element \(x \neq 0\) such that \(\operatorname{ad} x\) is nilpotent. (Indeed, we then have \(\mathfrak{g}' \neq 0\) .)
\end{enumerate}

¶ 11) Let g be a Lie algebra, r its radical.

\begin{enumerate}
\def\labelenumi{\alph{enumi})}
\item
  Let \(\mathfrak{s}\) be a Levi subalgebra of \(\mathfrak{g}\) , \(\mathfrak{k}\) a Cartan subalgebra of \(\mathfrak{s}\) . Show that \(\mathfrak{k}\) is contained in a Cartan subalgebra \(\mathfrak{h}\) of \(\mathfrak{g}\) which is the sum of \(\mathfrak{k}\) and a subalgebra of \(\mathfrak{r}\) . (Use §2, Th. 2, Prop. 10 and Cor. 2 of Th. 1.)
\item
  Let \(\mathfrak{h}'\) be a Cartan subalgebra of \(\mathfrak{g}\) . Show that there exists a Levi subalgebra \(\mathfrak{s}'\) of \(\mathfrak{g}\) such that \(\mathfrak{h}'\) is the sum of a Cartan subalgebra of \(\mathfrak{s}'\) and a subalgebra of \(\mathfrak{r}\) . (The subalgebras \(\mathfrak{s}, \mathfrak{k}, \mathfrak{h}\) in \(a\) ) can be chosen so that \(\mathfrak{h} + \mathfrak{r} = \mathfrak{h}' + \mathfrak{r}\) . Put \(\mathfrak{a} = \mathfrak{h} + \mathfrak{r}\) , which is solvable. By Th. 3, there exists \(x \in \mathcal{C}^{\infty}(\mathfrak{a})\) such that \(e^{\mathrm{ad}_{\mathfrak{a}}x}\mathfrak{h} = \mathfrak{h}'\) . Then \(e^{\mathrm{ad}_{\mathfrak{g}}x}\) is a special automorphism of \(\mathfrak{g}\) which transforms \(\mathfrak{s}\) into the required Levi subalgebra.)
\item
  Let \(\mathfrak{s}\) be a Levi subalgebra of \(\mathfrak{g}\) . Let \(\mathfrak{h}\) be a Cartan subalgebra of \(\mathfrak{g}\) which is the sum of a Cartan subalgebra \(\mathfrak{k}\) of \(\mathfrak{s}\) and a subalgebra \(\mathfrak{l}\) of \(\mathfrak{r}\) . Let \(\mathfrak{c}\) be the commutant of \(\mathfrak{k}\) in \(\mathfrak{r}\) . Show that \(\mathfrak{l}\) is a Cartan subalgebra of \(\mathfrak{c}\) . (For \(x \in \mathfrak{k}\) , \(\operatorname{ad}_{\mathfrak{g}} x\) is semi-simple, but \(\operatorname{ad}_{\mathfrak{h}} x\) is nilpotent, so \([\mathfrak{k}, \mathfrak{h}] = 0\) . If \(y \in \mathfrak{c}\) is such that \([y, \mathfrak{l}] \subset \mathfrak{l}\) , then \([y, \mathfrak{h}] \subset \mathfrak{h}\) so \(y \in \mathfrak{h} \cap \mathfrak{r} = \mathfrak{l}\) .)
\item
  Let \(\mathfrak{s}\) be a Levi subalgebra of \(\mathfrak{g}\) , \(\mathfrak{k}\) a Cartan subalgebra of \(\mathfrak{s}\) , \(\mathfrak{c}\) the commutant of \(\mathfrak{k}\) in \(\mathfrak{r}\) , \(\mathfrak{l}\) a Cartan subalgebra of \(\mathfrak{c}\) . Then \(\mathfrak{h} = \mathfrak{k} + \mathfrak{l}\) is a Cartan subalgebra of \(\mathfrak{g}\) . (Let \(x = y + z\) ( \(y \in \mathfrak{s}, z \in \mathfrak{r}\) ) be an element of the normalizer \(\mathfrak{u}\) of \(\mathfrak{h}\) in \(\mathfrak{g}\) . Show that \([y, \mathfrak{k}] \subset \mathfrak{k}\) , so \(y \in \mathfrak{k}\) and \(z \in \mathfrak{u}\) . Then show that \([z, \mathfrak{k}] \subset \mathfrak{h} \cap \mathfrak{r} \subset \mathfrak{c}\) , so \([\mathfrak{k}, [\mathfrak{k}, z]] = 0\) , and hence that \([\mathfrak{k}, z] = 0\) and \(z \in \mathfrak{c}\) . Finally, \([z, \mathfrak{l}] \subset \mathfrak{l}\) hence \(z \in \mathfrak{l}\) and \(x \in \mathfrak{h}\) .)
\item
  Let \(\mathfrak{s},\mathfrak{k},\mathfrak{c}\) be as in \(d\) ), and \(\mathfrak{q}=[\mathfrak{g},\mathfrak{r}]\) the nilpotent radical of \(\mathfrak{g}\) . Let \(x\in\mathfrak{q}\) and \(u\) the special automorphism \(e^{\mathrm{ad}x}\) . If \(u(\mathfrak{k})\subset\mathfrak{k}+\mathfrak{c}\) , then \(x\in\mathfrak{c}\) . (Consider the adjoint representation \(\rho\) of \(\mathfrak{s}\) on \(\mathfrak{q}\) , and let \(\mathfrak{q}_{i}\) be a complement of \(\mathcal{C}^{i+1}\mathfrak{q}\) in \(\mathcal{C}^{i}\mathfrak{q}\) stable under \(\rho\) ; let \(\rho_{i}\) be the subrepresentation of \(\rho\) defined by \(\mathfrak{q}_{i}\) . Let \(\sigma_{i}=\rho_{i}|\mathfrak{k}\) . Let \(\mathfrak{q}_{i}^{\prime}\) be the commutant of \(\mathfrak{k}\) in \(\mathfrak{q}_{i}\) and \(\mathfrak{q}_{i}^{\prime\prime}\) a complement of \(\mathfrak{q}_{i}^{\prime}\) in \(\mathfrak{q}_{i}\) stable under \(\sigma_{i}\) . Let \(x=x_{1}^{\prime}+x_{1}^{\prime\prime}+\cdots+x_{n}^{\prime}+x_{n}^{\prime\prime}\) with \(x_{i}^{\prime}\in\mathfrak{q}_{i}^{\prime},x_{i}^{\prime\prime}\in\mathfrak{q}_{i}^{\prime\prime}\) . Arguing by contradiction, assume that the \(x_{i}^{\prime\prime}\) are not all zero and \(x_{1}^{\prime\prime}=\cdots=x_{p-1}^{\prime\prime}=0\) , \(x_{p}^{\prime\prime}\neq0\) , for example. If \(h\in\mathfrak{k}\) , \(u(h)=h+[x_{p}^{\prime\prime},h]+y\) with \(y\in\mathcal{C}^{p+1}\mathfrak{q}\) . Since \(u(h)\in\mathfrak{k}+\mathfrak{c}\) , this gives \([x_{p}^{\prime\prime},h]+y\in\mathfrak{q}_{p}^{\prime}+\mathfrak{q}_{p+1}^{\prime}+\cdots+\mathfrak{q}_{n}^{\prime}\) , hence \([h,x_{p}^{\prime\prime}]\in\mathfrak{q}_{p}^{\prime}\) , so \([h,x_{p}^{\prime\prime}]=0\) . Then \(x_{p}^{\prime\prime}=0\) , a contradiction.)
\end{enumerate}

\(f\) ) Let \(\mathfrak{h}\) be a Cartan subalgebra of \(\mathfrak{g}\) . Then \(\mathfrak{h}\) can be expressed uniquely as the sum of \(\mathfrak{h} \cap \mathfrak{r}\) and a Cartan subalgebra of a Levi subalgebra of \(\mathfrak{g}\) . (For the uniqueness, use \(e\) ) and Th. 5 of Chap. I, §6, no. 8.) The Levi subalgebra in question is not unique in general.

\(g\) ) Let \(\mathfrak{h}\) be a Cartan subalgebra of \(\mathfrak{g}\) , and \(\mathfrak{t} = \mathfrak{g}^0 (\mathfrak{h} \cap \mathfrak{r})\) . Then \(\mathfrak{h}\) is a Cartan subalgebra of \(\mathfrak{t}\) . We have \(\mathfrak{g} = \mathfrak{t} + \mathfrak{r}\) (use a Fitting decomposition for the adjoint representation of \(\mathfrak{h} \cap \mathfrak{r}\) on \(\mathfrak{g}\) ). The algebra \(\mathfrak{t} \cap \mathfrak{r}\) is the radical of \(\mathfrak{t}\) and is nilpotent (use Exerc. 5 of §2).

\begin{enumerate}
\def\labelenumi{\alph{enumi})}
\setcounter{enumi}{7}
\tightlist
\item
  Assume that k is algebraically closed. Let h be a Cartan subalgebra of g. There exists a Levi subalgebra s of g such that, for all \(\lambda \in h^{*}\) ,
\end{enumerate}

\[
\mathfrak {g} ^ {\lambda} (\mathfrak {h}) = \left(\mathfrak {g} ^ {\lambda} (\mathfrak {h}) \cap \mathfrak {s}\right) + \left(\mathfrak {g} ^ {\lambda} (\mathfrak {h}) \cap \mathfrak {r}\right).
\]

(With the notations of g), take for s a Levi subalgebra of t such that \(\mathfrak{h} = (\mathfrak{h} \cap \mathfrak{s}) + (\mathfrak{h} \cap \mathfrak{r})\) ; this exists by b).) \(^{7}\)

¶ 12) a) Let g be a solvable Lie algebra, and G a finite subgroup (resp. compact subgroup if k = R or C) of Aut(g). Show that there exists a Cartan subalgebra of g stable under G. (Argue by induction on dim g, and reduce to the case in which g is an extension of a nilpotent algebra g/n by a commutative ideal n which is a non-trivial simple g/n-module (cf.~proof of Th. 3). The Cartan subalgebras of g then form an affine space attached to n, on which G operates. Conclude by an argument with the barycentre.)

\begin{enumerate}
\def\labelenumi{\alph{enumi})}
\setcounter{enumi}{1}
\tightlist
\item
  Let g be a solvable Lie algebra and s a Lie subalgebra of \(\text{Der}(\mathfrak{g})\) . Assume that the s-module g is semi-simple. Show that there exists a Cartan subalgebra of g stable under s. (Same method.)
\end{enumerate}

¶ 13) Let g be a Lie algebra and G a finite subgroup of Aut(g). Assume that G is hyper-solvable (Algebra, Chap. I, §6, Exerc. 26). Show that there exists a Cartan subalgebra of g stable under G.

(Argue by induction on dim g. Using Exerc. 12, reduce to the case where g is semi-simple. If G ≠ \{1\}, choose a normal subgroup C of G that is cyclic of prime order (Algebra, loc. cit.). The subalgebra s consisting of the elements invariant under C is reductive in g (§1, no. 5), and distinct from g. By the induction hypothesis, s has a Cartan subalgebra a stable under G. We have s ≠ 0 (Chap. I, §4, Exerc. 21 c)), so a ≠ 0. The commutant z of a in g is distinct from g and stable under G; choose a Cartan subalgebra h of z stable under G, and show that h is a Cartan subalgebra of g, cf.~the Cor. to Prop. 3.)

Construct a finite group of automorphisms of \(\mathfrak{sl}(2,\mathbf{C})\) which is isomorphic to \(A_{4}\) (and hence solvable) and which does not leave any Cartan subalgebra stable.

\begin{enumerate}
\def\labelenumi{\arabic{enumi})}
\setcounter{enumi}{13}
\item
  *Show that every irreducible complex (resp. real) linear representation of a hyper-solvable finite group G is induced by a representation of degree 1 (resp. of degree 1 or 2) of a subgroup of G. (Apply Exerc. 13 to the Lie algebras \(\mathfrak{gl}(n,\mathbf{C})\) and \(\mathfrak{gl}(n,\mathbf{R}).\) )
\item
  Assume that k is algebraically closed. Let g be a Lie algebra, h a Cartan subalgebra of g, and A a subset of g. Assume that A is dense in g (in the Zariski topology) and stable under \(\operatorname{Aut}_{e}(\mathfrak{g})\) . Show that \(A \cap h\) is dense in h. (Let X be the closure of \(A \cap h\) , and \(U = h - X\) . Assume that \(U \neq \varnothing\) . With the notations in Lemma 2, the image under F of \(U \times \mathfrak{g}^{\lambda_{1}}(\mathfrak{h}) \times \cdots \times \mathfrak{g}^{\lambda_{p}}(\mathfrak{h})\) contains a non-empty open subset of g. Since this image is contained in g-A, this contradicts the fact that A is dense in g.)
\item
  Let V be a finite dimensional k-vector space, and g a Lie subalgebra of \(\mathfrak{gl}(V)\) . We are going to show that the following three properties are equivalent: (i) The Cartan subalgebras of g are commutative and consist of semi-simple elements.
\end{enumerate}

\begin{enumerate}
\def\labelenumi{(\roman{enumi})}
\setcounter{enumi}{1}
\item
  Every regular element of \(\mathfrak{g}\) is semi-simple.
\item
  The semi-simple elements of g are dense in g in the Zariski topology.
\end{enumerate}

\begin{enumerate}
\def\labelenumi{\alph{enumi})}
\item
  Show that (i) \(\Longrightarrow\) (ii) \(\Longrightarrow\) (iii).
\item
  Let A be the set of semi-simple elements of \(\mathfrak{g}\) . Show that A is stable under \(\mathrm{Aut}_e(\mathfrak{g})\) .
\item
  Show that (iii) \(\Longrightarrow\) (i). (Reduce (App. I, Exerc. 1) to the case where k is algebraically closed. If h is a Cartan subalgebra of g, show, by using Exerc. 15, that \(A \cap h\) is dense in h; since \([x, y] = 0\) if \(x \in A \cap h, y \in h\) , deduce that h is commutative, from which (i) is immediate.)
\item
  Assume that k is R, C or a non-discrete complete ultrametric field of characteristic zero. Give g the topology defined by that of k. Show that properties (i), (ii), (iii) are equivalent to the following:
\end{enumerate}

\begin{enumerate}
\def\labelenumi{(\roman{enumi})}
\setcounter{enumi}{3}
\tightlist
\item
  The semi-simple elements of \(\mathfrak{g}\) are dense in \(\mathfrak{g}\) .
\end{enumerate}

(Show that (iv) \(\Longrightarrow\) (iii) and (ii) \(\Longrightarrow\) (iv), cf.~App. I, Exerc. 4.)

\begin{enumerate}
\def\labelenumi{\arabic{enumi})}
\setcounter{enumi}{16}
\tightlist
\item
  Assume that \(k\) is algebraically closed. Let \(\mathfrak{g}\) be a Lie algebra, \(\mathfrak{h}\) a Cartan subalgebra, and A a subset of the centre of \(\mathfrak{h}\) . Denote by \(\mathrm{E}_{\mathfrak{g}}\) the subgroup of \(\operatorname{Aut}(\mathfrak{g})\) denoted by E in no. 2. Show that, if \(s\) is an element of \(\operatorname{Aut}(\mathfrak{g})\) such that \(s\mathrm{A} = \mathrm{A}\) , there exists \(t \in \mathrm{E}_{\mathfrak{g}}\) such that \(t\mathfrak{h} = \mathfrak{h}\) and \(t|\mathrm{A} = \mathrm{Id}_{\mathrm{A}}\) ; in particular, \(ts|_{\mathrm{A}} = s|_{\mathrm{A}}\) . (Let \(\mathfrak{a}\) be the commutant of A in \(\mathfrak{g}\) ; since \(\mathfrak{h}\) and \(sh\) are Cartan subalgebras of \(\mathfrak{a}\) , there exists \(\theta \in \mathrm{E}_{\mathfrak{a}}\) such that \(\theta(s\mathfrak{h}) = \mathfrak{h}\) ; choose \(t\) from the elements of \(\mathrm{E}_{\mathfrak{g}}\) that extend \(\theta\) .)
\end{enumerate}

¶ 18) Let g be a Lie algebra, h a Cartan subalgebra of g and Ug (resp. Uh) the enveloping algebra of g (resp. h). A linear form \(\varphi\) on Ug is said to be central if it vanishes on {[}Ug, Ug{]}, i.e.~if \(\varphi(a.b) = \varphi(b.a)\) for all \(a, b \in Ug\) .

\(a)\) Let \(x, y \in \mathfrak{g}\) . Assume that there exists \(s \in \operatorname{Aut}_e(\mathfrak{g})\) such that \(s(x) = y\) . Show that

\[
\varphi (x ^ {n}) = \varphi (y ^ {n})
\]

for all \(n \in \mathbf{N}\) and for every central linear form \(\varphi\) on \(\mathrm{Ug}\) .

\begin{enumerate}
\def\labelenumi{\alph{enumi})}
\setcounter{enumi}{1}
\item
  Let \(\varphi\) be a central linear form on Ug whose restriction to Uh vanishes. Show that \(\varphi = 0\) . (We can assume that k is algebraically closed. Deduce from a) that we then have \(\varphi(x^{n}) = 0\) for all \(n \in N\) and all regular \(x \in g\) ; use a density argument to remove the assumption of regularity.)
\item
  Show that \(Ug = [Ug, Ug] + Uh\) .
\end{enumerate}

\(d\) ) Let V be a semi-simple \(\mathfrak{g}\) -module. Show that V is semi-simple as an \(\mathfrak{h}\) -module. In particular, \(V^{\lambda}(\mathfrak{h}) = V_{\lambda}(\mathfrak{h})\) for all \(\lambda \in \mathfrak{h}^{*}\) .

\begin{enumerate}
\def\labelenumi{\alph{enumi})}
\setcounter{enumi}{4}
\tightlist
\item
  Let \(V'\) be a semi-simple g-module. Assume that V and \(V'\) are isomorphic as h-modules. Show that they are isomorphic as g-modules. (If \(a \in Uh\) , remark that \(\operatorname{Tr}(a_{\mathrm{V}}) = \operatorname{Tr}(a_{\mathrm{V}'})\) ). Deduce, by using b) or c), that \(\operatorname{Tr}(x_{\mathrm{V}}) = \operatorname{Tr}(x_{\mathrm{V}'})\) for all \(x \in Ug\) , and conclude by using Algebra, Chap. VIII.)
\end{enumerate}

If k is algebraically closed, the assumption that ``V and V' are h-isomorphic'' is equivalent to saying that \(\dim\mathrm{V}_{\lambda}(\mathfrak{h})=\dim\mathrm{V}_{\lambda}^{\prime}(\mathfrak{h})\) for all \(\lambda\in\mathfrak{h}^{*}\) .

\exercisesection{\S~4}

The notations and assumptions are those of nos. 1, 2, 3 of §4.

\begin{enumerate}
\def\labelenumi{\arabic{enumi})}
\tightlist
\item
  Take \(\mathrm{G} = \mathbf{GL}_{n}(k), n \geq 0.\)
\end{enumerate}

\begin{enumerate}
\def\labelenumi{\alph{enumi})}
\item
  Show that \(r_{\mathrm{Ad}}^0 (g) = n\) for all \(g \in \mathrm{G}\) .
\item
  Show that an element \(g \in G\) is regular if and only if its characteristic polynomial \(\mathrm{P}_{g}(\mathrm{T}) = \det(\mathrm{T} - g)\) is separable; this is equivalent to saying that the discriminant (Algebra, Chap. IV, §1, no. 10) of \(\mathrm{P}_{g}(\mathrm{T})\) is \(\neq 0\) .
\end{enumerate}

\begin{enumerate}
\def\labelenumi{\arabic{enumi})}
\setcounter{enumi}{1}
\item
  Construct a Lie group G such that the function \(r_{Ad}^{0}\) is not constant. (Take g abelian \(\neq 0\) and Ad non-trivial.)
\item
  Let \((\rho_{i})_{i\in I}\) be a countable family of analytic linear representations of G. Prove that the elements of G which are regular for all the \(\rho_{i}\) constitute a dense subset of G. Construct an example showing that the countability assumption cannot be omitted.
\item
  Assume that \(k = \mathbf{C}\) and that \(\mathbf{G}\) is connected. Prove the equivalence of the following properties:
\end{enumerate}

\begin{enumerate}
\def\labelenumi{(\roman{enumi})}
\item
  G is nilpotent.
\item
  Every element \(\neq 1\) of \(G\) is regular.
\end{enumerate}

(Show first that (ii) implies

\begin{enumerate}
\def\labelenumi{(\roman{enumi})}
\setcounter{enumi}{1}
\tightlist
\item
  \(^{\prime}\) Every element \(\neq 0\) of g is regular.
\end{enumerate}

Remark next that, if \(g \neq 0\) , then g contains elements \(x \neq 0\) such that ad x is nilpotent, cf.~§3, Exerc. 10. Deduce that g is nilpotent, hence (i).

\exercisesection{\S~5}

\begin{enumerate}
\def\labelenumi{\arabic{enumi})}
\item
  Show that the solvable Lie algebra considered in Chap. I, §5, Exerc. 6 is not isomorphic to any decomposable Lie algebra.
\item
  Let u (resp. v) be a non-zero semi-simple (resp. nilpotent) endomorphism of V. Then the map \(\lambda u \mapsto \lambda v (\lambda \in k)\) is an isomorphism from g = ku to \(g' = kv\) which does not take semi-simple elements to semi-simple elements.
\item
  Let \(u\) be an endomorphism of \(V\) that is neither semi-simple nor nilpotent. Then \(\mathfrak{g} = ku\) is non-decomposable, but \(\operatorname{ad}_{\mathfrak{g}}\mathfrak{g}\) is decomposable.
\end{enumerate}

¶ 4) Let g be a decomposable Lie subalgebra of gl(V). Let q be a Lie subalgebra of g whose identity representation is semi-simple. There exists \(a \in \operatorname{Aut}_{e}(\mathfrak{g})\) such that \(a(\mathfrak{q})\) is contained in the subalgebra m of Prop. 7. (Imitate the proof of Prop. 7 (ii).)

¶ 5) Let g be a Lie subalgebra of gl(V). Then g is said to be algebraic if, for all \(x \in \mathfrak{g}\) , the replicas of \(x\) (Chap. I, §5, Exerc. 14) belong to g. Such an algebra is decomposable.

\begin{enumerate}
\def\labelenumi{\alph{enumi})}
\tightlist
\item
  Denote by \(a(\mathfrak{g})\) the smallest algebraic subalgebra of \(\mathfrak{gl}(\mathrm{V})\) containing \(\mathfrak{g}\) . Then
\end{enumerate}

\[
a (\mathfrak {g}) \supset e (\mathfrak {g}) \supset \mathfrak {g}.
\]

Give an example where \(a(\mathfrak{g})\) and \(e(\mathfrak{g})\) are distinct (take V of dimension 2 and \(\mathfrak{g}\) of dimension 1).

\begin{enumerate}
\def\labelenumi{\alph{enumi})}
\setcounter{enumi}{1}
\item
  Show that, if \(\mathfrak{n}\) is an ideal of \(\mathfrak{g}\) , \(\mathfrak{n}\) and \(a(\mathfrak{n})\) are ideals of \(a(\mathfrak{g})\) , and \([a(\mathfrak{g}), a(\mathfrak{n})] = [\mathfrak{g}, \mathfrak{n}]\) (imitate the proof of Prop. 4). Deduce that \(\mathcal{D}^i a(\mathfrak{g}) = \mathcal{D}^i \mathfrak{g}\) for \(i \geq 1\) and that \(\mathcal{C}^i a(\mathfrak{g}) = \mathcal{C}^i \mathfrak{g}\) for \(i \geq 2\) .
\item
  Show that every Lie algebra consisting of nilpotent elements is algebraic. \(^{8}\)
\end{enumerate}

¶ 6) Let g be a semi-simple subalgebra of \(\mathfrak{gl}(V)\) , \(\mathbf{T}(V) = \bigoplus_{n=0}^{\infty} \mathbf{T}^{n}(V)\) the tensor algebra of V, \(\mathbf{T}(V)^{\mathfrak{g}}\) the set of elements of \(\mathbf{T}(V)\) invariant under g (cf.~Chap. III, App.) and \(\overline{g}\) the set of \(u \in \mathfrak{gl}(V)\) such that u.x = 0 for all \(x \in \mathbf{T}(V)^{\mathfrak{g}}\) . We are going to show that \(\overline{g} = g\) .

\begin{enumerate}
\def\labelenumi{\alph{enumi})}
\item
  Show that the representation of g on the dual V* of V is isomorphic to its representation on \(\bigwedge^{p-1}V\) , where \(p = \dim V\) (use the fact that g is contained in \(\mathfrak{sl}(V)\) ). Deduce that every element of \(\mathbf{T}_{n,m} = \mathbf{T}^{n}(\mathrm{V}) \otimes \mathbf{T}^{m}(\mathrm{V}^{*})\) invariant under g is invariant under \(\overline{g}\) , and that \(\overline{g}\) is algebraic (Exerc. 5).
\item
  Let W be a vector subspace of one of the \(T_{n,m}\) . Assume that W is stable under g. Show that W is stable under \(\overline{g}\) (if \(e_{1},\ldots,e_{r}\) is a basis of W, remark that \(e_{1}\wedge\cdots\wedge e_{r}\) is invariant under g, and hence under \(\overline{g}\) ).
\end{enumerate}

Deduce that g is an ideal of \(\overline{g}\) , and that \(\overline{g}/g\) is commutative (cf.~proof of Prop. 4). We have \(\overline{g} = g \times c\) , where c is the centre of \(\overline{g}\) .

\begin{enumerate}
\def\labelenumi{\alph{enumi})}
\setcounter{enumi}{2}
\item
  Let R be the associative subalgebra of \(\mathfrak{gl}(V)\) generated by 1 and g. Show that c is contained in the centre of R (remark that \(\overline{g}\) is contained in the bicommutant of R, which is equal to R). Deduce that the elements of c are semi-simple.
\item
  Let \(x \in c\) . Show that the replicas of x belong to c (Chap. I, §5, Exerc. 14). Show that \(\operatorname{Tr}(sx) = 0\) for all \(s \in g\) ; deduce that \(\operatorname{Tr}(sx) = 0\) for all \(s \in \overline{g}\) , and hence that x is nilpotent (loc. cit).
\item
  Show that \(\mathfrak{c} = 0\) and \(\overline{\mathfrak{g}} = \mathfrak{g}\) by combining \(c\) and \(d\) ).
\end{enumerate}

\begin{enumerate}
\def\labelenumi{\arabic{enumi})}
\setcounter{enumi}{6}
\tightlist
\item
  Let \(\mathfrak{g}\) be a Lie subalgebra of \(\mathfrak{gl}(\mathrm{V})\) . Let \(m, n\) be two integers \(\geq 0\) , W and \(\mathrm{W}'\) two vector subspaces of \(\mathbf{T}^m(\mathrm{V}) \otimes \mathbf{T}^n(\mathrm{V}^*)\) where \(\mathrm{V}^*\) is the dual of V. Assume that \(\mathrm{W}' \subset \mathrm{W}\) and that W and \(\mathrm{W}'\) are stable under the natural representation of \(\mathfrak{g}\) on \(\mathbf{T}^m(\mathrm{V}) \otimes \mathbf{T}^n(\mathrm{V}^*)\) . Show that W and \(\mathrm{W}'\) are then stable under \(e(\mathfrak{g})\) . If \(\pi\) denotes the representation of \(e(\mathfrak{g})\) on \(\mathrm{W/W}'\) thus obtained, show that \(\pi e(\mathfrak{g})\) is the decomposable envelope of \(\pi(\mathfrak{g})\) (use Th. 1).
\end{enumerate}

Deduce that \(\operatorname{ad} e(\mathfrak{g})\) is the decomposable envelope of \(\operatorname{ad} \mathfrak{g}\) in \(\mathfrak{gl}(\mathfrak{g})\) .

\begin{enumerate}
\def\labelenumi{\arabic{enumi})}
\setcounter{enumi}{7}
\tightlist
\item
  Let g be a Lie subalgebra of \(\mathfrak{gl}(V)\) and h a Cartan subalgebra of g.
\end{enumerate}

\begin{enumerate}
\def\labelenumi{\alph{enumi})}
\tightlist
\item
  Show that \(e(\mathfrak{g}) = e(\mathfrak{h}) + \mathcal{D}\mathfrak{g} = e(\mathfrak{h}) + \mathfrak{g}\) .
\end{enumerate}

(Remark that \(e(\mathfrak{h}) + \mathcal{D}\mathfrak{g}\) is decomposable (Cor. 1 of Th. 1), contains \(\mathfrak{g} = \mathfrak{h} + \mathcal{D}\mathfrak{g}\) , and is contained in \(e(\mathfrak{g})\) ; thus, it is \(e(\mathfrak{g})\) .)

\begin{enumerate}
\def\labelenumi{\alph{enumi})}
\setcounter{enumi}{1}
\item
  We have \(e(\mathfrak{h}) \cap \mathfrak{g} = \mathfrak{h}\) (remark that \(e(\mathfrak{h}) \cap \mathfrak{g}\) is nilpotent).
\item
  Let \(x\) be an element of the normalizer of \(e(\mathfrak{h})\) in \(e(\mathfrak{g})\) . Show that \(x \in e(\mathfrak{h})\) . (Write \(x = y + z\) with \(y \in e(\mathfrak{h}), z \in \mathfrak{g}\) , cf.~a); remark that \([z, \mathfrak{h}] \subset e(\mathfrak{h}) \cap \mathfrak{g} = \mathfrak{h}\) , hence \(z \in \mathfrak{h}\) .)
\end{enumerate}

\(d\) ) Show that \(e(\mathfrak{h})\) is a Cartan subalgebra of \(e(\mathfrak{g})\) .

\begin{enumerate}
\def\labelenumi{\arabic{enumi})}
\setcounter{enumi}{8}
\tightlist
\item
  Let \(\mathfrak{g}\) be a Lie subalgebra of \(\mathfrak{gl}(\mathrm{V})\) . Show that conditions (i), (ii), (iii) of Exerc. 16 of §3 are equivalent to:
\end{enumerate}

\begin{enumerate}
\def\labelenumi{(\alph{enumi})}
\setcounter{enumi}{21}
\tightlist
\item
  \(\mathfrak{g}\) is decomposable and has the same rank as \(\mathfrak{g} / \mathfrak{n}_{\mathrm{V}}(\mathfrak{g})\) .
\end{enumerate}

(If \(\mathfrak{h}\) is a Cartan subalgebra of \(\mathfrak{g}\) , the condition `` \(\mathfrak{g}\) has the same rank as \(\mathfrak{g}/\mathfrak{n}_{\mathrm{V}}(\mathfrak{g})\) '' is equivalent to saying that \(\mathfrak{h} \cap \mathfrak{n}_{\mathrm{V}}(\mathfrak{g}) = 0\) , i.e.~that \(\mathfrak{h}\) has no nilpotent element \(\neq 0\) (cf.~§2, Exerc. 14). Deduce the equivalence of (i) and (v).)

\begin{enumerate}
\def\labelenumi{\arabic{enumi})}
\setcounter{enumi}{9}
\tightlist
\item
  Let \(k'\) be an extension of \(k\) and \(\mathfrak{g}'\) a \(k'\) -Lie subalgebra of
\end{enumerate}

\[
\mathfrak {g l} (\mathrm{V} \otimes_ {k} k ^ {\prime}) = \mathfrak {g l} (\mathrm{V}) \otimes_ {k} k ^ {\prime}.
\]

\begin{enumerate}
\def\labelenumi{\alph{enumi})}
\item
  Show that there exists a smallest Lie subalgebra g of \(\mathfrak{gl}(V)\) such that \(g\otimes_{k}k'\) contains \(g'\) .
\item
  Assume that \(k'\) is algebraically closed and denote by G the group of k-automorphisms of \(k'\) ; this group operates in a natural way on \(V \otimes_{k} k'\) . Show that \(g \otimes_{k} k'\) is the Lie subalgebra generated by the conjugates of \(g'\) under G (use the fact that k is the field of G-invariants in \(k'\) ).
\item
  Show that \(\mathfrak{g}\) is decomposable if \(\mathfrak{g}'\) is decomposable. (Reduce to the case where \(k'\) is algebraically closed, and use \(b\) ) as well as Cor. 1 and 3 of Th. 1.)
\end{enumerate}

¶ 11) Exceptionally, we assume in this exercise that k is a perfect field of characteristic p \textgreater{} 0.

Let g be a Lie p-algebra (Chap. I, §1, Exerc. 20). If \(x \in g\) , denote by \(\langle x \rangle\) the smallest Lie p-subalgebra containing x. It is commutative and generated as a k-vector space by the \(x^{p^{i}}\) where \(i = 0, 1, \ldots\) . Then x is said to be nilpotent (resp. semi-simple) if the p-map of \(\langle x \rangle\) is nilpotent (resp. bijective).

\begin{enumerate}
\def\labelenumi{\alph{enumi})}
\item
  Show that x can be decomposed uniquely in the form \(x = s + n\) , with \(s, n \in \langle x \rangle\) , s semi-simple and n nilpotent (apply Exerc. 23 of Chap. I, §1). If f is a p-homomorphism from g to \(\mathfrak{gl}(\mathrm{V})\) , \(f(s)\) and \(f(n)\) are the semi-simple and nilpotent components of the endomorphism \(f(x)\) ; this applies in particular to f = ad.
\item
  A subalgebra of g is said to be decomposable if it contains the semi-simple and nilpotent components of its elements. Show that, if b and c are vector subspaces of g such that \(b \subset c\) , the set of \(x \in g\) such that \([x, c] \subset b\) is decomposable (same proof as Prop. 3); in particular, every Cartan subalgebra of g is decomposable.
\item
  Let t be a commutative subalgebra of g consisting of semi-simple elements, and maximal with this property. Let h be the commutant of t in g. Let \(x \in h\) and let \(x = s + n\) be its canonical decomposition; since h is decomposable (cf.~b)), \(s, n \in h\) . Show that the subalgebra generated by t and s is commutative and consists of semi-simple elements, hence coincides with t. Deduce that \(ad_{h}x = ad_{h}n\) is nilpotent, and hence that h is nilpotent. Since \(\mathfrak{h} = \mathfrak{g}^{0}(\mathfrak{h})\) , h is a Cartan subalgebra of g (§2, Prop. 4).
\end{enumerate}

In particular, every Lie p-algebra over a finite field has a Cartan subalgebra. \(^{9}\)

\exercisesection{Appendix~I}

Denote by V a finite dimensional vector space over \(k\) .

\begin{enumerate}
\def\labelenumi{\arabic{enumi})}
\item
  Let \(k'\) be an extension of k, and let \(\mathrm{V}_{(k')} = \mathrm{V} \otimes_k k'\) . Show that the Zariski topology on \(\mathrm{V}_{(k')}\) induces the Zariski topology on V, and that V is dense in \(\mathrm{V}_{(k')}\) .
\item
  Assume that \(\mathrm{V}\) is the product of two vector spaces \(\mathrm{V}_1\) and \(\mathrm{V}_2\) .
\end{enumerate}

\begin{enumerate}
\def\labelenumi{\alph{enumi})}
\item
  The Zariski topology on V is finer than the product topology of the Zariski topologies on \(V_{1}\) and \(V_{2}\) ; it is strictly finer if \(V_{1} \neq 0\) and \(V_{2} \neq 0\) .
\item
  If \(A_{1}\) (resp. \(A_{2}\) ) is a subset of \(V_{1}\) (resp. \(V_{2}\) ), the closure of \(A_{1} \times A_{2}\) is the product of the closures of \(A_{1}\) and \(A_{2}\) .
\end{enumerate}

\begin{enumerate}
\def\labelenumi{\arabic{enumi})}
\setcounter{enumi}{2}
\tightlist
\item
  Assume that \(k\) is algebraically closed. Let \(A\) and \(B\) be two closed subsets of \(V\) , and \(\mathfrak{a}\) (resp. \(\mathfrak{b}\) ) the set of \(f \in A_V\) which vanish on \(A\) (resp. \(B\) ). Prove the equivalence of the following properties:
\end{enumerate}

\begin{enumerate}
\def\labelenumi{(\roman{enumi})}
\item
  A ∩ B = ∅.
\item
  \(\mathfrak{a} + \mathfrak{b} = \mathrm{A_V}\) .
\item
  There exists a polynomial function \(f\) on \(V\) which is equal to 1 on \(A\) and to 0 on \(B\) .
\end{enumerate}

(Use Hilbert's theorem of zeros (Commutative Algebra, Chap. V, §3, no. 3) to prove that (i) \(\Longrightarrow\) (ii).)

\begin{enumerate}
\def\labelenumi{\arabic{enumi})}
\setcounter{enumi}{3}
\tightlist
\item
  Assume that \(k\) is a non-discrete complete valued field. Denote by \(\mathcal{T}\) (resp. \(\mathcal{Z}\) ) the Banach space (resp. Zariski) topology on V.
\end{enumerate}

\begin{enumerate}
\def\labelenumi{\alph{enumi})}
\item
  Show that \(\mathcal{T}\) is finer than \(\mathcal{L}\) (and strictly finer if \(V \neq 0\) ).
\item
  Show that every non-empty \(\mathcal{Z}\) -open subset of V is \(\mathcal{T}\) -dense.
\end{enumerate}

\exercisesection{Appendix~II}

¶ 1) Let X be a locally connected topological space, \(\mathcal{C}(\mathrm{X})\) the space of continuous real-valued functions on X, and \(d\) an integer \(\geq 0\) . Let \(\mathrm{F} \in \mathcal{C}(\mathrm{X})[\mathrm{T}]\) be a monic polynomial of degree \(d\) with coefficients in \(\mathcal{C}(\mathrm{X})\) :

\[
\mathrm{F} = \mathrm{T} ^ {d} + \mathrm{T} ^ {d - 1} f _ {1} + \dots + f _ {d}, \quad f _ {i} \in \mathcal {C} (\mathrm{X}).
\]

Identify F with a function on \(R \times X\) by putting

\[
\mathrm{F} (t, x) = t ^ {d} + t ^ {d - 1} f _ {1} (x) + \dots + f _ {d} (x) \text {if} t \in \mathbf {R}, x \in \mathrm{X}.
\]

Let \(\Delta \in \mathcal{C}(\mathrm{X})\) be the discriminant of the polynomial F (Algebra, Chap. IV, §1, no. 10).

If U is an open subset of X, denote by \(Z_{U}\) the set of \((t,x)\) , with \(t \in R\) and \(x \in U\) , such that \(F(t,x) = 0\) ; this is a closed subset of \(R \times U\) .

\begin{enumerate}
\def\labelenumi{\alph{enumi})}
\item
  Show that the projection \(\mathrm{pr}_2: \mathrm{Z}_{\mathrm{U}} \to \mathrm{U}\) is proper (General Topology, Chap. I, §10).
\item
  Assume that U is connected and that \(\Delta(x) \neq 0\) for all \(x \in U\) . Show that \(Z_{U} \to U\) is a covering of U (General Topology, Chap. XI) of degree \(\leq d\) , and that the number of connected components of \(R \times U - Z_{U}\) is \(\leq d + 1\) .
\item
  Let \(X'\) be the set of points of X at which \(\Delta\) is \(\neq 0\) . Assume that \(X'\) is dense in X. Denote by \(\mathcal{A}(\text{resp. } \mathcal{B})\) the set of connected components of \(X'\) (resp. of \(R \times X - Z_{X}\) ). Show that
\end{enumerate}

\[
\operatorname{Card} (\mathscr {B}) \leq (d + 1) \operatorname{Card} (\mathscr {A}) \quad (\text { use   } b)).
\]

\begin{enumerate}
\def\labelenumi{\alph{enumi})}
\setcounter{enumi}{3}
\tightlist
\item
  Assume that X is connected, and that \(d \geq 1\) . Show that
\end{enumerate}

\[
\operatorname{Card} (\mathcal {B}) \leq 1 + d \operatorname{Card} (\mathcal {A}).
\]

¶ 2) Let V be a real vector space of finite dimension n, and F a polynomial function on V of degree d.~Let \(V'\) be the set of points of V at which \(F \neq 0\) . Show that the number of connected components of \(V'\) is finite and bounded by a constant depending only on n and d.~(Argue by induction on n.~Reduce to the case in which F has no multiple factors, and show that V can be decomposed as \(R \times X\) in such a way that the results of Exerc. 1 are applicable to F.) \(^{10}\)

\chapter{Split Semi-simple Lie Algebras}

In this chapter, k denotes a (commutative) field of characteristic 0. Unless otherwise stated, by a ``vector space'', we mean a ``vector space over k''; similarly for ``Lie algebra'', etc.

\section{\texorpdfstring{THE LIE ALGEBRA \(\mathfrak{sl}(2,k)\) AND ITS REPRESENTATIONS}{THE LIE ALGEBRA \textbackslash mathfrak\{sl\}(2,k) AND ITS REPRESENTATIONS}}

\subsection{\texorpdfstring{CANONICAL BASIS OF \(\mathfrak{sl}(2,k)\)}{CANONICAL BASIS OF \textbackslash mathfrak\{sl\}(2,k)}}

Lemma 1. Let \(\mathbf{A}\) be an associative algebra over \(k\) , \(H\) and \(X\) elements of \(\mathbf{A}\) such that \([H, X] = 2X\) .

\begin{enumerate}
\def\labelenumi{(\roman{enumi})}
\item
  \([H,X^n ] = 2nX^n\) for any integer \(n\geq 0\)
\item
  If \(Z\) is an element of \(A\) such that \([Z, X] = H\) , then, for any integer \(n > 0\) ,
\end{enumerate}

\[
[ Z, X ^ {n} ] = n X ^ {n - 1} (H + n - 1) = n (H - n + 1) X ^ {n - 1}.
\]

The map \(T \mapsto [H, T]\) from \(A\) to \(A\) is a derivation, which implies (i). With the assumptions in (ii),

\[
\begin{array}{l} [ Z, X ^ {n} ] = \sum_ {i + j = n - 1} X ^ {i} H X ^ {j} \\ \qquad = \sum_ {i + j = n - 1} (X ^ {i} X ^ {j} H + X ^ {i} 2 j X ^ {j}) \\ \qquad = n X ^ {n - 1} H + 2 X ^ {n - 1} \frac {n (n - 1)}{2} \\ \qquad = n X ^ {n - 1} (H + n - 1). \end{array}
\]

On the other hand, \(X^{n-1}(H+n-1)=(H-n+1)X^{n-1}\) by (i). Q.E.D.

Recall that we denote by \(\mathfrak{sl}(2,k)\) the Lie algebra consisting of the square matrices of order 2, trace zero, and with entries in k. This Lie algebra is simple of dimension 3 (Chap. I, §6, no. 7, Example). The canonical basis of \(\mathfrak{sl}(2,k)\) is the basis \((X_{+},X_{-},H)\) , where

\[
X _ {+} = \left( \begin{array}{c c} 0 & 1 \\ 0 & 0 \end{array} \right) \quad X _ {-} = \left( \begin{array}{c c} 0 & 0 \\ - 1 & 0 \end{array} \right) \quad H = \left( \begin{array}{c c} 1 & 0 \\ 0 & - 1 \end{array} \right).
\]

We have

\[
[ H, X _ {+} ] = 2 X _ {+} \quad [ H, X _ {-} ] = - 2 X _ {-} \quad [ X _ {+}, X _ {-} ] = - H.\tag{1}
\]

Since the identity representation of \(\mathfrak{sl}(2,k)\) is injective, H is a semi-simple element of \(\mathfrak{sl}(2,k)\) and \(X_{+}, X_{-}\) are nilpotent elements of \(\mathfrak{sl}(2,k)\) (Chap. I, §6, no. 3, Th. 3). By Chap. VII, §2, no. 1, Example 4, kH is a Cartan subalgebra of \(\mathfrak{sl}(2,k)\) . The map \(U \mapsto -^{t}U\) is an involutive automorphism of the Lie algebra \(\mathfrak{sl}(2,k)\) , called the canonical involution of \(\mathfrak{sl}(2,k)\) ; it transforms \((X_{+}, X_{-}, H)\) into \((X_{-}, X_{+}, -H)\) .

Lemma 2. In the enveloping algebra of \(\mathfrak{sl}(2,k)\) ,

\[
[ H, X _ {+} ^ {n} ] = 2 n X _ {+} ^ {n} \quad [ H, X _ {-} ^ {n} ] = - 2 n X _ {-} ^ {n}
\]

for any integer \(n \geq 0\) , and

\[
[ X _ {-}, X _ {+} ^ {n} ] = n X _ {+} ^ {n - 1} (H + n - 1) = n (H - n + 1) X _ {+} ^ {n - 1}
\]

\[
[ X _ {+}, X _ {-} ^ {n} ] = n X _ {-} ^ {n - 1} (- H + n - 1) = n (- H - n + 1) X _ {-} ^ {n - 1}
\]

if n \textgreater{} 0.

The first and third relations follow from Lemma 1. The others can be deduced from them by using the canonical involution of \(\mathfrak{sl}(2,k)\) .

\subsection{\texorpdfstring{PRIMITIVE ELEMENTS OF \(\mathfrak{sl}(2,k)\) -MODULES}{PRIMITIVE ELEMENTS OF \textbackslash mathfrak\{sl\}(2,k) -MODULES}}

Let E be an \(\mathfrak{sl}(2,k)\) -module. If \(A \in \mathfrak{sl}(2,k)\) and \(x \in E\) , we shall often write Ax instead of \(A_{E}x\) . Let \(\lambda \in k\) . If Hx = \(\lambda x\) we say, by abuse of language, that x is an element of E of weight \(\lambda\) , or that \(\lambda\) is the weight of x. If E is finite dimensional, \(H_{E}\) is semi-simple, so the set of elements of weight \(\lambda\) is the primary subspace of E relative to \(H_{E}\) and \(\lambda\) (cf.~Chap. VII, §1, no. 1).

Lemma 3. If \(x\) is an element of weight \(\lambda\) , then \(X_{+}x\) is an element of weight \(\lambda + 2\) and \(X_{-}x\) is an element of weight \(\lambda - 2\) .

Indeed, \(HX_{+}x = [H,X_{+}]x + X_{+}Hx = 2X_{+}x + X_{+}\lambda x = (\lambda +2)X_{+}x\) , and similarly \(HX_{-}x = (\lambda -2)X_{-}x\) (cf.~also Chap. VII, §1, no. 3, Prop. 10 (ii)).

DEFINITION 1. Let E be an \(\mathfrak{sl}(2,k)\) -module. An element of E is said to be primitive if it is a non-zero eigenvector of \(H_{E}\) and belongs to the kernel of \(X_{+E}\) .

A non-zero element e of E is primitive if and only if ke is stable under the operation of \(kH + kX_{+}\) ; this follows for example from Lemma 3.

Examples The element \(X_{+}\) is primitive of weight 2 for the adjoint representation of \(\mathfrak{sl}(2,k)\) . The element \((1,0)\) of \(k^{2}\) is primitive of weight 1 for the identity representation of \(\mathfrak{sl}(2,k)\) on \(k^{2}\) .

Lemma 4. Let \(\mathbf{E}\) be a non-zero finite dimensional \(\mathfrak{sl}(2,k)\) -module. Then \(\mathbf{E}\) has primitive elements.

Since \(X_{+}\) is a nilpotent element of \(\mathfrak{sl}(2,k)\) , \(X_{+E}\) is nilpotent. Assume that \(X_{+E}^{m-1} \neq 0\) and \(X_{+E}^{m} = 0\) . By Lemma 2,

\[
m (H _ {\mathrm{E}} - m + 1) X _ {+ \mathrm{E}} ^ {m - 1} = [ X _ {- \mathrm{E}}, X _ {+ \mathrm{E}} ^ {m} ] = 0,
\]

and hence the elements of \(X_{+}^{m - 1}(\mathrm{E}) - \{0\}\) are primitive.

PROPOSITION 1. Let E be an \(\mathfrak{sl}(2,k)\) -module, and \(e\) a primitive element of E of weight \(\lambda\) . Put \(e_n = \frac{(-1)^n}{n} X_-^n e\) for \(n \geq 0\) , and \(e_{-1} = 0\) . Then

\[
\left\{ \begin{array}{l l} H e _ {n} & = (\lambda - 2 n) e _ {n} \\ X _ {-} e _ {n} & = - (n + 1) e _ {n + 1} \\ X _ {+} e _ {n} & = (\lambda - n + 1) e _ {n - 1}. \end{array} \right.\tag{2}
\]

The first formula follows from Lemma 3, and the second from the definition of the \(e_n\) . We prove the third by induction on \(n\) . It is satisfied for \(n = 0\) since \(e_{-1} = 0\) . If \(n > 0\) ,

\[
\begin{array}{r l} & n X _ {+} e _ {n} = - X _ {+} X _ {-} e _ {n - 1} = - [ X _ {+}, X _ {-} ] e _ {n - 1} - X _ {-} X _ {+} e _ {n - 1} \\ & \qquad = H e _ {n - 1} - X _ {-} (\lambda - n + 2) e _ {n - 2} \\ & \qquad = (\lambda - 2 n + 2 + (n - 1) (\lambda - n + 2)) e _ {n - 1} \\ & \qquad = n (\lambda - n + 1) e _ {n - 1}. \end{array}
\]

COROLLARY. The submodule of \(\mathbf{E}\) generated by \(e\) is the vector subspace generated by the \(e_n\) .

This follows from the formulas (2).

The integers \(n \geq 0\) such that \(e_{n} \neq 0\) constitute an interval in N, and the corresponding elements \(e_{n}\) form a basis over k of the submodule generated by e (indeed, they are linearly independent because they are non-zero elements of distinct weights). This basis will be said to be associated to the primitive element e.

PROPOSITION 2. If the submodule V of E generated by the primitive element e is finite dimensional, then:

\begin{enumerate}
\def\labelenumi{(\roman{enumi})}
\item
  the weight \(\lambda\) of e is integral and equal to \(\dim V - 1\) ;
\item
  \((e_0, e_1, \ldots, e_\lambda)\) is a basis of V, and \(e_n = 0\) for \(n > \lambda\) ;
\item
  the eigenvalues of \(H_{\mathrm{V}}\) are \(\lambda, \lambda - 2, \lambda - 4, \ldots, -\lambda\) ; they are all of multiplicity 1;
\item
  every primitive element of V is proportional to e;
\item
  the commutant of the module \(\mathrm{V}\) is reduced to the scalars; in particular, \(\mathrm{V}\) is absolutely simple.
\end{enumerate}

Let \(m\) be the largest integer such that \(e_m \neq 0\) . Then \(0 = X_{+}e_{m+1} = (\lambda - m)e_m\) , so \(\lambda = m\) ; since \((e_0, e_1, \ldots, e_m)\) is a basis of \(V\) , this proves (i) and (ii). Assertion (iii) follows from the equality \(He_n = (\lambda - 2n)e_n\) . We have \(X_{+}e_n \neq 0\) for \(1 \leq n \leq m\) , hence (iv). Let \(c\) be an element of the commutant of the module \(V\) . Then \(Hc(e) = cH(e) = \lambda c(e)\) , so there exists \(\mu \in k\) such that \(c(e) = \mu e\) ; then

\[
c X _ {-} ^ {q} e = X _ {-} ^ {q} c e = \mu X _ {-} ^ {q} e
\]

for all \(q \geq 0\) , so \(c = \mu.1\) , proving (v).

COROLLARY. Let E be a finite dimensional \(\mathfrak{sl}(2,k)\) -module.

\begin{enumerate}
\def\labelenumi{(\roman{enumi})}
\item
  The endomorphism \(H_{E}\) is diagonalizable and its eigenvalues are rational integers.
\item
  For any \(p \in Z\) , let \(E_{p}\) be the eigenspace of \(H_{E}\) corresponding to the eigenvalue p.~Let i be an integer \(\geq 0\) . The map \(X_{-E}^{i}|E_{p}:E_{p} \to E_{p-2i}\) is injective for \(i \leq p\) , bijective for i = p, and surjective for \(i \geq p\) . The map \(X_{+E}^{i}|E_{-p}:E_{-p} \to E_{-p+2i}\) is injective for \(i \leq p\) , bijective for i = p, and surjective for \(i \geq p\) .
\item
  The length of \(\mathrm{E}\) is equal to \(\dim \operatorname{Ker} X_{+\mathrm{E}}\) and to \(\dim \operatorname{Ker} X_{-\mathrm{E}}\) .
\item
  Let \(E'\) (resp. \(E''\) ) be the sum of the \(E_{p}\) for p even (resp. odd). Then \(E'\) (resp. \(E''\) ) is the sum of the simple submodules of E of odd (resp. even) dimension; and \(E = E' \oplus E''\) . The length of \(E'\) is \(\dim E_{0}\) , and that of \(E''\) is \(\dim E_{1}\) .
\end{enumerate}

\[
\text {(v)} \operatorname{Ker} X _ {+ \mathrm{E}} \cap \operatorname{Im} X _ {+ \mathrm{E}} \subset \sum_ {p > 0} \mathrm{E} _ {p} \text {and} \operatorname{Ker} X _ {- \mathrm{E}} \cap \operatorname{Im} X _ {- \mathrm{E}} \subset \sum_ {p <   0} \mathrm{E} _ {p}.
\]

If E is simple, E is generated by a primitive element (Lemma 4), and it suffices to apply Propositions 1 and 2. The general case follows since every finite dimensional \(\mathfrak{sl}(2,k)\) -module is semi-simple.

\subsection{THE SIMPLE MODULES V(m)}

Let \((u,v)\) be the canonical basis of \(k^2\) . For the identity representation of \(\mathfrak{sl}(2,k)\) ,

\[
X _ {+} u = 0 \quad H u = u \quad X _ {-} u = - v
\]

\[
X _ {+} v = u \quad H v = - v \quad X _ {-} v = 0.
\]

Consider the symmetric algebra \(\mathbf{S}(k^2)\) of \(k^2\) (Algebra, Chap. III, §6, no. 1, Def. 1). The elements of \(\mathfrak{sl}(2,k)\) extend uniquely to derivations of \(\mathbf{S}(k^2)\) , giving \(\mathbf{S}(k^2)\) the structure of an \(\mathfrak{sl}(2,k)\) -module (Chap. I, §3, no. 2). Let \(\mathrm{V}(m)\) be the set of homogeneous elements of \(\mathbf{S}(k^2)\) of degree \(m\) . Then \(\mathrm{V}(m)\) is an \(\mathfrak{sl}(2,k)\) -submodule of \(\mathbf{S}(k^2)\) of dimension \(m + 1\) , the \(m\) th symmetric power of \(\mathrm{V}(1) = k^2\) (Chap. III, Appendix). If \(m,n\) are integers such that \(0 \leq n \leq m\) , put

\[
e _ {n} ^ {(m)} = \binom {m} {n} u ^ {m - n} v ^ {n} \in \mathrm{V} (m).
\]

PROPOSITION 3. For any integer \(m \geq 0\) , \(V(m)\) is an absolutely simple \(\mathfrak{sl}(2, k)\) -module. In this module, \(e_0^{(m)} = u^m\) is primitive of weight \(m\) .

We have \(X_{+}u^{m}=0\) and \(Hu^{m}=mu^{m}\) , so \(u^{m}\) is primitive of weight m. The submodule of \(V(m)\) generated by \(u^{m}\) is of dimension \(m+1\) (Prop. 2(i)) and so is equal to \(V(m)\) . By Prop. 2(v), \(V(m)\) is absolutely simple.

THEOREM 1. Every simple \(\mathfrak{sl}(2,k)\) -module of finite dimension \(n\) is isomorphic to \(\mathrm{V}(n - 1)\) . Every finite dimensional \(\mathfrak{sl}(2,k)\) -module is a direct sum of submodules isomorphic to the modules \(\mathrm{V}(m)\) .

This follows from Lemma 4 and Prop. 1, 2 and 3.

Remarks. 1) The adjoint representation of \(\mathfrak{sl}(2,k)\) defines on \(\mathfrak{sl}(2,k)\) the structure of a simple \(\mathfrak{sl}(2,k)\) -module. This module is isomorphic to V(2) by an isomorphism that takes \(u^2\) to \(X_+\) , \(2uv\) to \(-H\) , and \(v^2\) to \(X_-\) .

\begin{enumerate}
\def\labelenumi{\arabic{enumi})}
\setcounter{enumi}{1}
\tightlist
\item
  For \(n \geq 0\) and m \textgreater{} n,
\end{enumerate}

\[
X _ {-} e _ {n} ^ {(m)} = - (m - n) \binom{m}{n} u ^ {m - n - 1} v ^ {n + 1} = - (n + 1) e _ {n + 1} ^ {(m)}.
\]

Hence, \((e_0^{(m)}, e_1^{(m)}, \ldots, e_m^{(m)})\) is the basis of \(V(m)\) associated to the primitive element \(e_0^{(m)}\) .

\begin{enumerate}
\def\labelenumi{\arabic{enumi})}
\setcounter{enumi}{2}
\tightlist
\item
  Let \(\Phi\) be the bilinear form on \(V(m)\) such that
\end{enumerate}

\[
\begin{array}{c} \varPhi (e _ {n} ^ {(m)}, e _ {n ^ {\prime}} ^ {(m)}) = 0 \text {if} n + n ^ {\prime} \neq m \\ \varPhi (e _ {n} ^ {(m)}, e _ {m - n} ^ {(m)}) = (- 1) ^ {n} \binom {m} {n}. \end{array}
\]

If \(x = au + bv\) and \(y = cu + dv\) , then \(\Phi(x^{m}, y^{m}) = (ad - bc)^{m}\) . It is now easy to check that \(\Phi\) is invariant, and that \(\Phi\) is symmetric for m even, and alternating for m odd.

PROPOSITION 4. Let E be a finite dimensional \(\mathfrak{sl}(2,k)\) -module, \(m\) an integer \(\geq 0\) , \(\mathrm{P}_m\) the set of primitive elements of weight \(m\) . Let L be the vector space of homomorphisms from the \(\mathfrak{sl}(2,k)\) -module \(\mathrm{V}(m)\) to the \(\mathfrak{sl}(2,k)\) -module E. The map \(f \mapsto f(u^m)\) from L to E is linear, injective, and its image is \(\mathrm{P}_m \cup \{0\}\) .

This map is clearly linear, and it is injective because \(u^{m}\) generates the \(\mathfrak{sl}(2,k)\) -module \(\mathrm{V}(m)\) . If \(f \in L\) ,

\[
X _ {+} (f (u ^ {m})) = f (X _ {+} u ^ {m}) = 0, \quad H (f (u ^ {m})) = f (H u ^ {m}) = m f (u ^ {m})
\]

so \(f(u^{m}) \in \mathrm{P}_{m} \cup \{0\}\) . Let \(e \in \mathrm{P}_{m}\) , and \(V\) the submodule of \(E\) generated by \(e\) . By Prop. 1, there exists an isomorphism from the module \(V(m)\) to the module \(V\) that takes \(u^{m}\) to \(e\) . Then \(L(u^{m}) = P_{m} \cup \{0\}\) .

COROLLARY. The isotypical component of E of type V(m) has length

\(\dim (\mathrm{P}_m\cup \{0\} .\)

\subsection{LINEAR REPRESENTATIONS OF THE GROUP SL(2, k)}

Recall (Algebra, Chap. III, §8, no. 9) that we denote by \(\mathbf{SL}(2,k)\) the group of square matrices of order 2 with coefficients in \(k\) whose determinant is equal to 1. If \(x\in \mathfrak{sl}(2,k)\) is nilpotent, then \(x^{2} = 0\) (Algebra, Chap. VII, §5, Cor. 3 of Prop. 5) and \(e^x = 1 + x\in \mathbf{SL}(2,k)\) . If E is a finite dimensional vector space and \(\rho\) is a linear representation of \(\mathfrak{sl}(2,k)\) on E, then \(\rho (x)\) is nilpotent and so \(e^{\rho (x)}\) is defined (Chap. I, §6, no. 3).

DEFINITION 2. Let E be a finite dimensional vector space, and \(\rho\) (resp. \(\pi\) ) a linear representation of \(\mathfrak{sl}(2,k)\) (resp. \(\mathbf{SL}(2,k)\) ) on E. Then \(\rho\) and \(\pi\) are said to be compatible if, for every nilpotent element x of \(\mathfrak{sl}(2,k)\) , \(\pi(e^{x}) = e^{\rho(x)}\) .

In other words, \(\rho\) and \(\pi\) are compatible if, for every nilpotent element x of \(\mathfrak{sl}(2,k)\) , the restriction of \(\rho\) to kx is compatible with the restriction of \(\pi\) to the group \(1 + kx\) (Chap. VII, §3, no. 1).

If \(\rho\) and \(\pi\) are compatible, so are the dual representations, the mth tensor powers, and the mth symmetric powers of \(\rho\) and \(\pi\) , respectively (Chap. VII, §5, no. 4, Lemma 1 (i) and (ii)). Similarly for the representations induced by \(\rho\) and \(\pi\) on a vector subspace stable under \(\rho\) and \(\pi\) (loc. cit.).

In particular, the representation \(\rho_{m}\) of \(\mathfrak{sl}(2,k)\) on \(V(m)\) (no. 3) is compatible with the mth symmetric power \(\pi_{m}\) of the identity representation \(\pi_{1}\) of \(\mathbf{SL}(2,k)\) . Putting \(e_{n}^{(m)} = \binom{m}{n} u^{m-n} v^{n}\) as above, we have

\[
\pi_ {m} (s) e _ {n} ^ {(m)} = \binom {m} {n} (s u) ^ {m - n} (s v) ^ {n}\tag{3}
\]

for \(s \in \mathbf{SL}(2, k)\) and \(0 \leq n \leq m\) .

THEOREM 2. Let \(\rho\) be a linear representation of \(\mathfrak{sl}(2,k)\) on a finite dimensional vector space E.

\begin{enumerate}
\def\labelenumi{(\roman{enumi})}
\item
  There exists a unique linear representation \(\pi\) of \(\mathbf{SL}(2,k)\) on E that is compatible with \(\rho\) .
\item
  A vector subspace \(\mathrm{F}\) of \(\mathrm{E}\) is stable under \(\pi\) if and only if it is stable under \(\rho\) .
\item
  Let \(x \in \mathrm{E}\) . Then \(\pi(s)x = x\) for all \(s \in \mathbf{SL}(2,k)\) if and only if \(x\) is invariant under \(\rho\) (that is, \(\rho(a)x = 0\) for all \(a \in \mathfrak{sl}(2,k)\) ).
\end{enumerate}

The existence of \(\pi\) follows from the preceding and Th. 1. On the other hand, we know that the group \(\mathbf{SL}(2,k)\) is generated by the elements of the form

\[
e ^ {t X _ {+}} = \left( \begin{array}{c c} 1 & t \\ 0 & 1 \end{array} \right) \quad e ^ {- t X _ {-}} = \left( \begin{array}{c c} 1 & 0 \\ t & 1 \end{array} \right)
\]

where \(t \in k\) (Algebra, Chap. III, §8, no. 9, Prop. 17). This proves the uniqueness of \(\pi\) .

Assertions (ii) and (iii) follow from what we have said, together with Chap. VII, §3, no. 1, Lemma 1 (i). Q.E.D.

Every finite dimensional \(\mathfrak{sl}(2,k)\) -module therefore has a unique \(\mathbf{SL}(2,k)\) -module structure, which is said to be associated to its \(\mathfrak{sl}(2,k)\) -module structure.

Remark. When k is R or C or a complete non-discrete ultrametric field, \(\mathfrak{sl}(2,k)\) is the Lie algebra of \(\mathbf{SL}(2,k)\) . Let \(\rho\) and \(\pi\) be as in Th. 2. The homomorphism \(\pi\) is a homomorphism of Lie groups from \(\mathbf{SL}(2,k)\) to \(\mathbf{GL}(\mathrm{E})\) : this is clear when \(\mathrm{E} = \mathrm{V}(m)\) , and the general case follows, in view of Th. 1. By Chap. VII, §3, no. 1, \(\rho(X_{+}) = \mathrm{L}(\pi)(X_{+})\) , \(\rho(X_{-}) = \mathrm{L}(\pi)(X_{-})\) . Hence \(\rho = \mathrm{L}(\pi)\) (for a converse, see Exerc. 18).

PROPOSITION 5. Let E, F be finite dimensional \(\mathfrak{sl}(2,k)\) -modules, and let \(f \in \operatorname{Hom}_{k}(\operatorname{E},\operatorname{F})\) . The following conditions are equivalent:

\begin{enumerate}
\def\labelenumi{(\roman{enumi})}
\item
  \(f\) is a homomorphism of \(\mathfrak{sl}(2,k)\) -modules;
\item
  \(f\) is a homomorphism of \(\mathbf{SL}(2,k)\) -modules.
\end{enumerate}

Condition (i) means that \(f\) is an invariant element of the \(\mathfrak{sl}(2,k)\) -module \(\mathrm{Hom}_k(\mathrm{E},\mathrm{F})\) , and condition (ii) means that \(f\) is an invariant element of the \(\mathbf{SL}(2,k)\) -module \(\mathrm{Hom}_k(\mathrm{E},\mathrm{F})\) . Since these module structures are associated by Chap. VII, §5, no. 4, Lemma 1 (iii), the proposition follows from Th. 2 (iii).

DEFINITION 3. The adjoint representation of the group \(\mathbf{SL}(2,k)\) is the linear representation Ad of \(\mathbf{SL}(2,k)\) on \(\mathfrak{sl}(2,k)\) defined by

\[
\operatorname{Ad} (s). a = s a s ^ {- 1}
\]

for all \(a \in \mathfrak{sl}(2, k)\) and all \(s \in \mathbf{SL}(2, k)\) .

When k is R or C or a complete non-discrete ultrametric field, we recover Def. 7 of Chap. III, §3, no. 12 (cf.~loc. cit., Prop. 49).

By Chap. VII, §5, no. 4, Lemma 1 (i) and (ii), the adjoint representations of \(\mathfrak{sl}(2,k)\) and \(\mathbf{SL}(2,k)\) are compatible. By Chap. VII, §3, no. 1, Remark 2, \(\operatorname{Ad}(\mathbf{SL}(2,k)) = \operatorname{Aut}_e(\mathfrak{sl}(2,k))\) .

\subsection{SOME ELEMENTS OF SL(2, k)}

For any \(t\in k^{*}\) , put

\[
\begin{array}{l} \theta (t) = e ^ {t X _ {+}} e ^ {t ^ {- 1} X _ {-}} e ^ {t X _ {+}} \\ \qquad = \left( \begin{array}{c c} 1 & t \\ 0 & 1 \end{array} \right) \left( \begin{array}{c c} 1 & 0 \\ - t ^ {- 1} & 1 \end{array} \right) \left( \begin{array}{c c} 1 & t \\ 0 & 1 \end{array} \right) \\ \qquad = \left( \begin{array}{c c} 0 & t \\ - t ^ {- 1} & 0 \end{array} \right) \\ \qquad = e ^ {t ^ {- 1} X _ {-}} e ^ {t X _ {+}} e ^ {t ^ {- 1} X _ {-}}. \end{array}
\]

With the notations of no. 3,

\[
\theta (t) u = - t ^ {- 1} v \quad \theta (t) v = t u
\]

so

\[
\theta (t) e _ {n} ^ {(m)} = (- 1) ^ {m - n} t ^ {2 n - m} e _ {m - n} ^ {(m)}.\tag{4}
\]

Hence, the element \(\theta(t)^2 = \left( \begin{array}{cc} -1 & 0 \\ 0 & -1 \end{array} \right)\) operates by \((-1)^m\) on \(\mathrm{V}(m)\) . If \(\mathrm{E}\) is an odd-dimensional simple \(\mathfrak{sl}(2,k)\) -module, \(\theta(t)_{\mathrm{E}}\) is thus an involutive automorphism of the vector space \(\mathrm{E}\) . In particular, taking \(\mathrm{E}\) to be the adjoint representation:

\[
\theta (t) _ {\mathrm{E}} X _ {+} = t ^ {- 2} X _ {-} \quad \theta (t) _ {\mathrm{E}} X _ {-} = t ^ {2} X _ {+} \quad \theta (t) _ {\mathrm{E}} H = - H\tag{5}
\]

so that \(\theta(1)_{\mathrm{E}} = \theta(-1)_{\mathrm{E}}\) is the canonical involution of \(\mathfrak{sl}(2,k)\) .

For any \(t\in k^{*}\) , put

\[
h (t) = \left( \begin{array}{c c} t & 0 \\ 0 & t ^ {- 1} \end{array} \right) = \theta (t) \theta (- 1).
\]

Then \(h(t)u = tu\) , \(h(t)v = t^{-1}v\) , so

\[
h (t) e _ {n} ^ {(m)} = t ^ {m - 2 n} e _ {n} ^ {(m)}.\tag{6}
\]

PROPOSITION 6. Let E be a finite dimensional \(\mathfrak{sl}(2,k)\) -module, and \(t\in k^{*}\) . Let \(\mathrm{E}_p\) be the set of elements of E of weight \(p\) .

\begin{enumerate}
\def\labelenumi{(\roman{enumi})}
\item
  \(\theta(t)_{\mathrm{E}}|\mathrm{E}_{p}\) is a bijection from \(\mathrm{E}_p\) to \(\mathrm{E}_{-p}\) .
\item
  \(h(t)_{\mathrm{E}}|\mathrm{E}_p\) is the homothety with ratio \(t^p\) on \(\mathrm{E}_p\) .
\end{enumerate}

If \(\mathrm{E} = \mathrm{V}(n)\) , the proposition follows from formulas (4) and (6). The general case follows from Th. 1.

COROLLARY. Let \(E = E' \oplus E''\) be the decomposition of E defined in the Cor. of Prop. 2. The element \(\begin{pmatrix}-1 & 0 \\ 0 & -1\end{pmatrix}\) of \(\mathbf{SL}(2, k)\) operates by +1 on \(E'\) and by -1 on \(E''\) .

This follows from (ii), applied to \(t = -1\) .

\section{ROOT SYSTEM OF A SPLIT SEMI-SIMPLE LIE ALGEBRA}

\subsection{SPLIT SEMI-SIMPLE LIE ALGEBRAS}

DEFINITION 1. Let \(\mathfrak{g}\) be a semi-simple Lie algebra. A Cartan subalgebra \(\mathfrak{h}\) of \(\mathfrak{g}\) is called splitting if, for all \(x \in \mathfrak{h}\) , \(\operatorname{ad}_{\mathfrak{g}} x\) is triangularizable. A semi-simple Lie algebra is called splittable if it has a splitting Cartan subalgebra. A split semi-simple Lie algebra is a pair \((\mathfrak{g}, \mathfrak{h})\) where \(\mathfrak{g}\) is a semi-simple Lie algebra and \(\mathfrak{h}\) is a splitting Cartan subalgebra of \(\mathfrak{g}\) .

Remarks. 1) Let \(\mathfrak{g}\) be a semi-simple Lie algebra, \(\mathfrak{h}\) a Cartan subalgebra of \(\mathfrak{g}\) . For all \(x \in \mathfrak{h}\) , \(\operatorname{ad}_{\mathfrak{g}} x\) is semi-simple (Chap. VII, §2, no. 4, Th. 2). Thus, to say that \(\mathfrak{h}\) is splitting means that \(\operatorname{ad}_{\mathfrak{g}} x\) is diagonalizable for all \(x \in \mathfrak{h}\) .

\begin{enumerate}
\def\labelenumi{\arabic{enumi})}
\setcounter{enumi}{1}
\item
  If \(k\) is algebraically closed, every semi-simple Lie algebra \(\mathfrak{g}\) is splittable, and every Cartan subalgebra of \(\mathfrak{g}\) is splitting. When \(k\) is not algebraically closed, there exist non-splittable semi-simple Lie algebras (Exerc. 2 a)); moreover, if \(\mathfrak{g}\) is splittable, there may exist Cartan subalgebras of \(\mathfrak{g}\) that are not splitting (Exerc. 2 b)).
\item
  Let \(\mathfrak{g}\) be a semi-simple Lie algebra, \(\mathfrak{h}\) a Cartan subalgebra of \(\mathfrak{g}\) , and \(\rho\) a finite dimensional injective representation of \(\mathfrak{g}\) such that \(\rho(\mathfrak{h})\) is diagonalizable. Then \(\operatorname{ad}_{\mathfrak{g}}x\) is diagonalizable for all \(x \in \mathfrak{h}\) (Chap. VII, §2, no. 1, Example 2), so \(\mathfrak{h}\) is splitting.
\item
  We shall see (§3, no. 3, Cor. of Prop. 10) that if \(\mathfrak{h}\) , \(\mathfrak{h}'\) are splitting Cartan subalgebras of \(\mathfrak{g}\) , there exists an elementary automorphism of \(\mathfrak{g}\) transforming \(\mathfrak{h}\) into \(\mathfrak{h}'\) .
\item
  Let \(\mathfrak{g}\) be a reductive Lie algebra. Then \(\mathfrak{g} = \mathfrak{c} \times \mathfrak{s}\) where \(\mathfrak{c}\) is the centre of \(\mathfrak{g}\) and \(\mathfrak{s} = \mathcal{D}\mathfrak{g}\) is semi-simple. The Cartan subalgebras of \(\mathfrak{g}\) are the subalgebras of the form \(\mathfrak{h} = \mathfrak{c} \times \mathfrak{h}'\) where \(\mathfrak{h}'\) is a Cartan subalgebra of \(\mathfrak{s}\) (Chap. VII, §2, no. 1, Prop. 2). Then \(\mathfrak{h}\) is called splitting if \(\mathfrak{h}'\) is splitting relative to \(\mathfrak{s}\) . This leads in an obvious way to the definition of splittable or split reductive algebras.
\end{enumerate}